\documentclass[a4paper,10pt]{article}
\usepackage{amsmath,amssymb}
\usepackage{fontspec}
\setmainfont{Liberation Serif}
\usepackage{graphicx}
\usepackage{eso-pic}
\usepackage{tikz}
\usepackage[most]{tcolorbox}
\usepackage{geometry}
\usepackage{fancyhdr}
\usepackage{lastpage}
\usepackage{booktabs}
\usepackage{listings}
\usepackage{enumitem}
\usepackage{microtype}

\geometry{
    a4paper,
    top=1.5cm,
    bottom=1.8cm,
    left=1.2cm,
    right=1.2cm
}

% Persistent Unremovable Canvas Watermark on Every Page (Shared XObject)
\newsavebox{\watermarkbox}
\sbox{\watermarkbox}{%
    \begin{tabular}{c}
        \includegraphics[width=8cm]{/home/johan/GATE/GATE_CS_PYQ_DB/assets/studiva_butterfly_icon_transparent.png} \\[-0.5em]
        {\fontsize{18}{22}\selectfont\bfseries\scshape www.studiva.co.in}
    \end{tabular}%
}
\AddToShipoutPictureBG{%
    \begin{tikzpicture}[remember picture, overlay]
        \node[opacity=0.12] at (current page.center) {\usebox{\watermarkbox}};
    \end{tikzpicture}%
}

% Header and Footer
\pagestyle{fancy}
\fancyhf{}
\fancyhead[L]{\small\scshape GATE 2025 Computer Science and Information Technology (CS1)}
\fancyhead[R]{\small\scshape Practice Question Bank with Solutions}
\fancyfoot[L]{\scriptsize\textsc{Copyright \copyright\ 2026} Studiva\textsuperscript{TM}}
\fancyfoot[C]{\small Page \thepage\ of \pageref{LastPage}}
\fancyfoot[R]{\raisebox{-0.2em}{\includegraphics[height=1.0em]{/home/johan/GATE/GATE_CS_PYQ_DB/assets/studiva_butterfly_icon_transparent.png}}\hspace{0.35em}\scriptsize\texttt{www.studiva.co.in}}
\renewcommand{\headrulewidth}{0.4pt}
\renewcommand{\footrulewidth}{0.4pt}

% Code Listing Style (Monochrome, clean)
\lstset{
    basicstyle=\footnotesize\ttfamily,
    numbers=none,
    frame=none,
    tabsize=2,
    breaklines=true,
    showstringspaces=false
}

\begin{document}

% ==================== PAGE 1: DEDICATED IMPRINT & COPYRIGHT PAGE ====================
\thispagestyle{empty}
\vspace*{0.5cm}
\begin{center}
    \includegraphics[width=2.5cm]{/home/johan/GATE/GATE_CS_PYQ_DB/assets/studiva_butterfly_icon_transparent.png}\\[0.6em]
    {\LARGE \textbf{GATE 2025 COMPUTER SCIENCE \& IT}}\\[0.3em]
    {\Large \textbf{PREVIOUS YEAR QUESTION BANK WITH SOLUTIONS}}\\[0.2em]
    {\normalsize \textsc{Session CS1 $\cdot$ Examination Year 2025}}\\[1.2em]
    \rule{0.6\textwidth}{0.6pt}
\end{center}

\vspace{1em}
\noindent\begin{tabular}{ll}
\textbf{Author / Compiler:} & Rupam Ghosh \\
\textbf{Publisher \& Platform:} & Studiva\textsuperscript{TM} (\texttt{www.studiva.co.in}) \\
\textbf{Official Inquiries:} & \texttt{contact@studiva.co.in} \\
\textbf{Edition:} & First Standard Edition (2026) \\
\textbf{Syllabus Reference:} & Official GATE 2026 Syllabus (Organizing Institute: IIT Guwahati) \\
\end{tabular}

\vspace{1.5em}
\begin{tcolorbox}[enhanced, standard jigsaw, opacityback=0, colframe=black!40, arc=1mm, boxrule=0.5pt, left=10pt, right=10pt, top=8pt, bottom=8pt]
\textbf{\large COPYRIGHT \& INTELLECTUAL PROPERTY NOTICE}\\[0.6em]
Copyright \copyright\ 2026 Rupam Ghosh \& Studiva\textsuperscript{TM}. All Rights Reserved.\\[0.5em]
This educational compilation, problem organization, taxonomy classification, LaTeX typesetting, hints, and explanatory solutions are the proprietary intellectual property of \textbf{Rupam Ghosh} and \textbf{Studiva\textsuperscript{TM}}.\\[0.6em]
\textbf{Terms of Use \& Attribution Rules:}
\begin{itemize}[leftmargin=*, nosep]
    \item \textbf{Personal Educational Use:} Permitted freely for individual students and educators preparing for competitive examinations.
    \item \textbf{Mandatory Attribution:} Any reproduction, republication, hosting, or sharing (digital or physical) must prominently credit \textbf{Rupam Ghosh \& Studiva\textsuperscript{TM}} and include an active, direct link to \texttt{https://www.studiva.co.in}.
    \item \textbf{Commercial Restrictions:} Commercial resale, sublicensing, course bundling, or inclusion in paid test series without prior written permission is strictly prohibited.
    \item \textbf{Anti-Tampering Notice:} Removing, altering, or obscuring copyright notices, page watermarks, or website links from this document is strictly prohibited.
\end{itemize}
\vspace{0.6em}
\textbf{Fair Use Disclaimer:}\\
The underlying examination questions originate from the Graduate Aptitude Test in Engineering (GATE), conducted by the Indian Institutes of Technology (IITs) and Indian Institute of Science (IISc). Question texts are reproduced for non-commercial educational study under Fair Use doctrines. No copyright is claimed over the original GATE examination question texts.
\end{tcolorbox}

\vspace{1.2em}
\begin{center}
\small\itshape
Practice Recommendation: Use an opaque sheet of paper or ruler to cover the right-hand column during mock attempts. Tick your chosen answers directly in the check boxes on the left, then slide the cover to check the solution and hint.
\end{center}

\newpage
% ==================== PAGE 2 ONWARDS: QUESTION BANK WITH SOLUTIONS ====================

\vspace{6pt}
\subsection*{General Aptitude}
\vspace{-4pt}
\noindent
\begin{tcolorbox}[
    blanker,
    sidebyside,
    sidebyside align=top,
    lefthand width=0.55\textwidth,
    top=4pt, bottom=4pt, left=0pt, right=0pt,
    middle=16pt,
    underlay={
        \draw[dashed, black!80, line width=0.8pt, dash pattern=on 4pt off 3pt]
            ([xshift=0.55\textwidth+8pt]frame.north west) --
            ([xshift=0.55\textwidth+8pt]frame.south west);
    }
]
\textbf{Q.1} \quad {\footnotesize\textsc{[MCQ $\cdot$ 1 Mark(s) $\cdot$ Neg: 0.33]}} \hfill {\tiny\texttt{2025\_CS1\_GA\_Q01}}
\vspace{2pt}\hrule height 0.3pt\vspace{4pt}

Ravi had \_\_\_\_\_\_ younger brother who taught at \_\_\_\_\_\_ university. He was widely regarded as \_\_\_\_\_\_ honorable man.
Select the option with the correct sequence of articles to fill in the blanks.

\par\vspace{6pt}\noindent
\begin{tabular}{@{}p{0.48\linewidth}@{\hspace{0.04\linewidth}}p{0.48\linewidth}@{}}
$\square$\ \textbf{(A)} a; a; an & $\square$\ \textbf{(B)} the; an; a \\
\addlinespace[3pt]
$\square$\ \textbf{(C)} a; an; a & $\square$\ \textbf{(D)} an; an; a
\end{tabular}

\tcblower

\textbf{Official Key:} \textbf{\large A} \hfill \textit{\small [Easy]}
\vspace{2pt}\hrule height 0.3pt\vspace{3pt}
\textbf{Topic:} Verbal Aptitude \\
\textbf{Subtopic:} \textit{Basic English grammar: tenses, articles, adjectives, prepositions, conjunctions, verb-noun agreement}
\vspace{2pt}\hrule height 0.3pt\vspace{3pt}
\textbf{Guide / Hint:} \\
{\small 'Younger' begins with a consonant sound (/j/), requiring the indefinite article 'a'. 'University' also begins with a consonant sound (/juː/), taking 'a'. 'Honorable' has a silent 'h' and begins with a vowel sound (/ɒ/), requiring 'an'. Hence, the sequence is a; a; an.}

\end{tcolorbox}
\vspace{6pt}\hrule height 0.35pt\vspace{6pt}
\noindent
\begin{tcolorbox}[
    blanker,
    sidebyside,
    sidebyside align=top,
    lefthand width=0.55\textwidth,
    top=4pt, bottom=4pt, left=0pt, right=0pt,
    middle=16pt,
    underlay={
        \draw[dashed, black!80, line width=0.8pt, dash pattern=on 4pt off 3pt]
            ([xshift=0.55\textwidth+8pt]frame.north west) --
            ([xshift=0.55\textwidth+8pt]frame.south west);
    }
]
\textbf{Q.2} \quad {\footnotesize\textsc{[MCQ $\cdot$ 1 Mark(s) $\cdot$ Neg: 0.33]}} \hfill {\tiny\texttt{2025\_CS1\_GA\_Q02}}
\vspace{2pt}\hrule height 0.3pt\vspace{4pt}

The CEO’s decision to downsize the workforce was considered myopic because it sacrificed long-term stability to accommodate short-term gains.
Select the most appropriate option that can replace the word ``myopic'' without changing the meaning of the sentence.

\par\vspace{6pt}\noindent
\begin{tabular}{@{}p{0.48\linewidth}@{\hspace{0.04\linewidth}}p{0.48\linewidth}@{}}
$\square$\ \textbf{(A)} visionary & $\square$\ \textbf{(B)} shortsighted \\
\addlinespace[3pt]
$\square$\ \textbf{(C)} progressive & $\square$\ \textbf{(D)} innovative
\end{tabular}

\tcblower

\textbf{Official Key:} \textbf{\large B} \hfill \textit{\small [Easy]}
\vspace{2pt}\hrule height 0.3pt\vspace{3pt}
\textbf{Topic:} Verbal Aptitude \\
\textbf{Subtopic:} \textit{Vocabulary: words, idioms, and phrases in context}
\vspace{2pt}\hrule height 0.3pt\vspace{3pt}
\textbf{Guide / Hint:} \\
{\small 'Myopic' literally means nearsighted, and metaphorically refers to lacking foresight or being shortsighted by focusing only on immediate considerations rather than long-term outcomes.}

\end{tcolorbox}
\vspace{6pt}\hrule height 0.35pt\vspace{6pt}
\noindent
\begin{tcolorbox}[
    blanker,
    sidebyside,
    sidebyside align=top,
    lefthand width=0.55\textwidth,
    top=4pt, bottom=4pt, left=0pt, right=0pt,
    middle=16pt,
    underlay={
        \draw[dashed, black!80, line width=0.8pt, dash pattern=on 4pt off 3pt]
            ([xshift=0.55\textwidth+8pt]frame.north west) --
            ([xshift=0.55\textwidth+8pt]frame.south west);
    }
]
\textbf{Q.3} \quad {\footnotesize\textsc{[MCQ $\cdot$ 1 Mark(s) $\cdot$ Neg: 0.33]}} \hfill {\tiny\texttt{2025\_CS1\_GA\_Q03}}
\vspace{2pt}\hrule height 0.3pt\vspace{4pt}

The average marks obtained by a class in an examination were calculated as 30.8. However, while checking the marks entered, the teacher found that the marks of one student were entered incorrectly as 24 instead of 42. After correcting the marks, the average becomes 31.4. How many students does the class have?

\par\vspace{6pt}\noindent
\begin{tabular}{@{}p{0.48\linewidth}@{\hspace{0.04\linewidth}}p{0.48\linewidth}@{}}
$\square$\ \textbf{(A)} 25 & $\square$\ \textbf{(B)} 28 \\
\addlinespace[3pt]
$\square$\ \textbf{(C)} 30 & $\square$\ \textbf{(D)} 32
\end{tabular}

\tcblower

\textbf{Official Key:} \textbf{\large C} \hfill \textit{\small [Easy]}
\vspace{2pt}\hrule height 0.3pt\vspace{3pt}
\textbf{Topic:} Quantitative Aptitude \\
\textbf{Subtopic:} \textit{Elementary statistics and probability}
\vspace{2pt}\hrule height 0.3pt\vspace{3pt}
\textbf{Guide / Hint:} \\
{\small Let $n$ be the number of students. The total increase in marks after correction is $42 - 24 = 18$. The increase in average is $31.4 - 30.8 = 0.6$. Therefore, $0.6 \times n = 18$, giving $n = \frac{18}{0.6} = 30$ students.}

\end{tcolorbox}
\vspace{6pt}\hrule height 0.35pt\vspace{6pt}
\noindent
\begin{tcolorbox}[
    blanker,
    sidebyside,
    sidebyside align=top,
    lefthand width=0.55\textwidth,
    top=4pt, bottom=4pt, left=0pt, right=0pt,
    middle=16pt,
    underlay={
        \draw[dashed, black!80, line width=0.8pt, dash pattern=on 4pt off 3pt]
            ([xshift=0.55\textwidth+8pt]frame.north west) --
            ([xshift=0.55\textwidth+8pt]frame.south west);
    }
]
\textbf{Q.4} \quad {\footnotesize\textsc{[MCQ $\cdot$ 1 Mark(s) $\cdot$ Neg: 0.33]}} \hfill {\tiny\texttt{2025\_CS1\_GA\_Q04}}
\vspace{2pt}\hrule height 0.3pt\vspace{4pt}

Consider the relationships among P, Q, R, S, and T:
\begin{itemize}
\item P is the brother of Q.
\item S is the daughter of Q.
\item T is the sister of S.
\item R is the mother of Q.
\end{itemize}
The following statements are made based on the relationships given above:
(1) R is the grandmother of S.
(2) P is the uncle of S and T.
(3) R has only one son.
(4) Q has only one daughter.
Which one of the following options is correct?

\par\vspace{6pt}\noindent
\begin{tabular}{@{}p{0.48\linewidth}@{\hspace{0.04\linewidth}}p{0.48\linewidth}@{}}
$\square$\ \textbf{(A)} Both (1) and (2) are true. & $\square$\ \textbf{(B)} Both (1) and (3) are true. \\
\addlinespace[3pt]
$\square$\ \textbf{(C)} Only (3) is true. & $\square$\ \textbf{(D)} Only (4) is true.
\end{tabular}

\tcblower

\textbf{Official Key:} \textbf{\large A} \hfill \textit{\small [Easy]}
\vspace{2pt}\hrule height 0.3pt\vspace{3pt}
\textbf{Topic:} Analytical Aptitude \\
\textbf{Subtopic:} \textit{Logic: deduction and induction}
\vspace{2pt}\hrule height 0.3pt\vspace{3pt}
\textbf{Guide / Hint:} \\
{\small Since R is the mother of Q and S is the daughter of Q, R is definitely the grandmother of S (statement 1 is true). P is the brother of Q, so P is the maternal/paternal uncle of Q's children S and T (statement 2 is true). The genders and number of other children of R and Q are not specified, so (3) and (4) cannot be established.}

\end{tcolorbox}
\vspace{6pt}\hrule height 0.35pt\vspace{6pt}
\noindent
\begin{tcolorbox}[
    blanker,
    sidebyside,
    sidebyside align=top,
    lefthand width=0.55\textwidth,
    top=4pt, bottom=4pt, left=0pt, right=0pt,
    middle=16pt,
    underlay={
        \draw[dashed, black!80, line width=0.8pt, dash pattern=on 4pt off 3pt]
            ([xshift=0.55\textwidth+8pt]frame.north west) --
            ([xshift=0.55\textwidth+8pt]frame.south west);
    }
]
\textbf{Q.6} \quad {\footnotesize\textsc{[MCQ $\cdot$ 2 Mark(s) $\cdot$ Neg: 0.66]}} \hfill {\tiny\texttt{2025\_CS1\_GA\_Q06}}
\vspace{2pt}\hrule height 0.3pt\vspace{4pt}

``I put the brown paper in my pocket along with the chalks, and possibly other things. I suppose every one must have reflected how primeval and how poetical are the things that one carries in one's pocket: the pocket-knife, for instance the type of all human tools, the infant of the sword. Once I planned to write a book of poems entirely about the things in my pocket. But I found it would be too long: and the age of the great epics is past.''
(From G.K. Chesterton's ``A Piece of Chalk'')
\par\medskip
Based only on the information provided in the above passage, which one of the following statements is true?

\par\vspace{6pt}\noindent
\begin{tabular}{@{}p{0.98\linewidth}@{}}
$\square$\ \textbf{(A)} The author of the passage carries a mirror in his pocket to reflect upon things. \\
\addlinespace[3pt]
$\square$\ \textbf{(B)} The author of the passage had decided to write a poem on epics. \\
\addlinespace[3pt]
$\square$\ \textbf{(C)} The pocket-knife is described as the infant of the sword. \\
\addlinespace[3pt]
$\square$\ \textbf{(D)} Epics are described as too inconvenient to write.
\end{tabular}

\tcblower

\textbf{Official Key:} \textbf{\large C} \hfill \textit{\small [Easy]}
\vspace{2pt}\hrule height 0.3pt\vspace{3pt}
\textbf{Topic:} Verbal Aptitude \\
\textbf{Subtopic:} \textit{Reading and comprehension}
\vspace{2pt}\hrule height 0.3pt\vspace{3pt}
\textbf{Guide / Hint:} \\
{\small The passage explicitly states: 'the pocket-knife, for instance the type of all human tools, the infant of the sword.' Statement C directly reflects this verbatim description.}

\end{tcolorbox}
\vspace{6pt}\hrule height 0.35pt\vspace{6pt}
\noindent
\begin{tcolorbox}[
    blanker,
    sidebyside,
    sidebyside align=top,
    lefthand width=0.55\textwidth,
    top=4pt, bottom=4pt, left=0pt, right=0pt,
    middle=16pt,
    underlay={
        \draw[dashed, black!80, line width=0.8pt, dash pattern=on 4pt off 3pt]
            ([xshift=0.55\textwidth+8pt]frame.north west) --
            ([xshift=0.55\textwidth+8pt]frame.south west);
    }
]
\textbf{Q.8} \quad {\footnotesize\textsc{[MCQ $\cdot$ 2 Mark(s) $\cdot$ Neg: 0.66]}} \hfill {\tiny\texttt{2025\_CS1\_GA\_Q08}}
\vspace{2pt}\hrule height 0.3pt\vspace{4pt}

A fair six-faced dice, with the faces labelled `1', `2', `3', `4', `5', and `6', is rolled thrice. What is the probability of rolling `6' exactly once?

\par\vspace{6pt}\noindent
\begin{tabular}{@{}p{0.48\linewidth}@{\hspace{0.04\linewidth}}p{0.48\linewidth}@{}}
$\square$\ \textbf{(A)} $\frac{75}{216}$ & $\square$\ \textbf{(B)} $\frac{1}{6}$ \\
\addlinespace[3pt]
$\square$\ \textbf{(C)} $\frac{1}{18}$ & $\square$\ \textbf{(D)} $\frac{25}{216}$
\end{tabular}

\tcblower

\textbf{Official Key:} \textbf{\large A} \hfill \textit{\small [Medium]}
\vspace{2pt}\hrule height 0.3pt\vspace{3pt}
\textbf{Topic:} Quantitative Aptitude \\
\textbf{Subtopic:} \textit{Elementary statistics and probability}
\vspace{2pt}\hrule height 0.3pt\vspace{3pt}
\textbf{Guide / Hint:} \\
{\small The probability of rolling a 6 is $p = \frac{1}{6}$, and not rolling a 6 is $q = \frac{5}{6}$. Using the binomial distribution for $n = 3$ trials with $k = 1$ success, $P(X = 1) = \binom{3}{1} \left(\frac{1}{6}\right)^1 \left(\frac{5}{6}\right)^2 = 3 \times \frac{1}{6} \times \frac{25}{36} = \frac{75}{216}$.}

\end{tcolorbox}
\vspace{6pt}\hrule height 0.35pt\vspace{6pt}
\noindent
\begin{tcolorbox}[
    blanker,
    sidebyside,
    sidebyside align=top,
    lefthand width=0.55\textwidth,
    top=4pt, bottom=4pt, left=0pt, right=0pt,
    middle=16pt,
    underlay={
        \draw[dashed, black!80, line width=0.8pt, dash pattern=on 4pt off 3pt]
            ([xshift=0.55\textwidth+8pt]frame.north west) --
            ([xshift=0.55\textwidth+8pt]frame.south west);
    }
]
\textbf{Q.10} \quad {\footnotesize\textsc{[MCQ $\cdot$ 2 Mark(s) $\cdot$ Neg: 0.66]}} \hfill {\tiny\texttt{2025\_CS1\_GA\_Q10}}
\vspace{2pt}\hrule height 0.3pt\vspace{4pt}

A shop has 4 distinct flavors of ice-cream. One can purchase any number of scoops of any flavor. The order in which the scoops are purchased is inconsequential. If one wants to purchase 3 scoops of ice-cream, in how many ways can one make that purchase?

\par\vspace{6pt}\noindent
\begin{tabular}{@{}p{0.48\linewidth}@{\hspace{0.04\linewidth}}p{0.48\linewidth}@{}}
$\square$\ \textbf{(A)} 4 & $\square$\ \textbf{(B)} 20 \\
\addlinespace[3pt]
$\square$\ \textbf{(C)} 24 & $\square$\ \textbf{(D)} 48
\end{tabular}

\tcblower

\textbf{Official Key:} \textbf{\large B} \hfill \textit{\small [Easy]}
\vspace{2pt}\hrule height 0.3pt\vspace{3pt}
\textbf{Topic:} Quantitative Aptitude \\
\textbf{Subtopic:} \textit{Numerical computation and estimation: ratios, percentages, powers, exponents and logarithms, permutations and combinations, series}
\vspace{2pt}\hrule height 0.3pt\vspace{3pt}
\textbf{Guide / Hint:} \\
{\small This is a multichoose problem (combination with repetition), where $r = 3$ scoops are chosen from $n = 4$ flavors. The number of ways is given by the stars and bars formula $\binom{n + r - 1}{r} = \binom{4 + 3 - 1}{3} = \binom{6}{3} = \frac{6 \times 5 \times 4}{3 \times 2 \times 1} = 20$.}

\end{tcolorbox}
\vspace{6pt}\hrule height 0.35pt\vspace{6pt}
\vspace{6pt}
\subsection*{Computer Science and Information Technology}
\vspace{-4pt}
\noindent
\begin{tcolorbox}[
    blanker,
    sidebyside,
    sidebyside align=top,
    lefthand width=0.55\textwidth,
    top=4pt, bottom=4pt, left=0pt, right=0pt,
    middle=16pt,
    underlay={
        \draw[dashed, black!80, line width=0.8pt, dash pattern=on 4pt off 3pt]
            ([xshift=0.55\textwidth+8pt]frame.north west) --
            ([xshift=0.55\textwidth+8pt]frame.south west);
    }
]
\textbf{Q.11} \quad {\footnotesize\textsc{[MCQ $\cdot$ 1 Mark(s) $\cdot$ Neg: 0.33]}} \hfill {\tiny\texttt{2025\_CS1\_CS\_Q11}}
\vspace{2pt}\hrule height 0.3pt\vspace{4pt}

Suppose a program is running on a non-pipelined single processor computer system. The computer is connected to an external device that can interrupt the processor asynchronously. The processor needs to execute the interrupt service routine (ISR) to serve this interrupt. The following steps (not necessarily in order) are taken by the processor when the interrupt arrives:
\begin{enumerate}[label=(\roman*)]
\item The processor saves the content of the program counter.
\item The program counter is loaded with the start address of the ISR.
\item The processor finishes the present instruction.
\end{enumerate}
Which ONE of the following is the CORRECT sequence of steps?

\par\vspace{6pt}\noindent
\begin{tabular}{@{}p{0.48\linewidth}@{\hspace{0.04\linewidth}}p{0.48\linewidth}@{}}
$\square$\ \textbf{(A)} (iii), (i), (ii) & $\square$\ \textbf{(B)} (i), (iii), (ii) \\
\addlinespace[3pt]
$\square$\ \textbf{(C)} (i), (ii), (iii) & $\square$\ \textbf{(D)} (iii), (ii), (i)
\end{tabular}

\tcblower

\textbf{Official Key:} \textbf{\large A} \hfill \textit{\small [Easy]}
\vspace{2pt}\hrule height 0.3pt\vspace{3pt}
\textbf{Topic:} Computer Organization and Architecture \\
\textbf{Subtopic:} \textit{I/O interface: interrupt and DMA mode}
\vspace{2pt}\hrule height 0.3pt\vspace{3pt}
\textbf{Guide / Hint:} \\
{\small When an asynchronous interrupt arrives, the processor first finishes the execution of the current instruction (iii). Next, it saves the current program counter (and processor state) onto the stack (i). Finally, it loads the starting address of the interrupt service routine into the program counter (ii).}

\end{tcolorbox}
\vspace{6pt}\hrule height 0.35pt\vspace{6pt}
\noindent
\begin{tcolorbox}[
    blanker,
    sidebyside,
    sidebyside align=top,
    lefthand width=0.55\textwidth,
    top=4pt, bottom=4pt, left=0pt, right=0pt,
    middle=16pt,
    underlay={
        \draw[dashed, black!80, line width=0.8pt, dash pattern=on 4pt off 3pt]
            ([xshift=0.55\textwidth+8pt]frame.north west) --
            ([xshift=0.55\textwidth+8pt]frame.south west);
    }
]
\textbf{Q.12} \quad {\footnotesize\textsc{[MCQ $\cdot$ 1 Mark(s) $\cdot$ Neg: 0.33]}} \hfill {\tiny\texttt{2025\_CS1\_CS\_Q12}}
\vspace{2pt}\hrule height 0.3pt\vspace{4pt}

Which ONE of the following statements is FALSE regarding the symbol table?

\par\vspace{6pt}\noindent
\begin{tabular}{@{}p{0.98\linewidth}@{}}
$\square$\ \textbf{(A)} Symbol table is responsible for keeping track of the scope of variables. \\
\addlinespace[3pt]
$\square$\ \textbf{(B)} Symbol table can be implemented using a binary search tree. \\
\addlinespace[3pt]
$\square$\ \textbf{(C)} Symbol table is not required after the parsing phase. \\
\addlinespace[3pt]
$\square$\ \textbf{(D)} Symbol table is created during the lexical analysis phase.
\end{tabular}

\tcblower

\textbf{Official Key:} \textbf{\large C} \hfill \textit{\small [Easy]}
\vspace{2pt}\hrule height 0.3pt\vspace{3pt}
\textbf{Topic:} Compiler Design \\
\textbf{Subtopic:} \textit{Lexical analysis, parsing, syntax-directed translation}
\vspace{2pt}\hrule height 0.3pt\vspace{3pt}
\textbf{Guide / Hint:} \\
{\small The symbol table is actively used throughout all phases of the compiler, including semantic analysis, intermediate code generation, code optimization, and target code generation for type checking, memory allocation, and address resolution. Thus, claiming it is not required after parsing is false.}

\end{tcolorbox}
\vspace{6pt}\hrule height 0.35pt\vspace{6pt}
\noindent
\begin{tcolorbox}[
    blanker,
    sidebyside,
    sidebyside align=top,
    lefthand width=0.55\textwidth,
    top=4pt, bottom=4pt, left=0pt, right=0pt,
    middle=16pt,
    underlay={
        \draw[dashed, black!80, line width=0.8pt, dash pattern=on 4pt off 3pt]
            ([xshift=0.55\textwidth+8pt]frame.north west) --
            ([xshift=0.55\textwidth+8pt]frame.south west);
    }
]
\textbf{Q.13} \quad {\footnotesize\textsc{[MCQ $\cdot$ 1 Mark(s) $\cdot$ Neg: 0.33]}} \hfill {\tiny\texttt{2025\_CS1\_CS\_Q13}}
\vspace{2pt}\hrule height 0.3pt\vspace{4pt}

Which ONE of the following techniques used in compiler code optimization uses live variable analysis?

\par\vspace{6pt}\noindent
\begin{tabular}{@{}p{0.98\linewidth}@{}}
$\square$\ \textbf{(A)} Run-time function call management \\
\addlinespace[3pt]
$\square$\ \textbf{(B)} Register assignment to variables \\
\addlinespace[3pt]
$\square$\ \textbf{(C)} Strength reduction \\
\addlinespace[3pt]
$\square$\ \textbf{(D)} Constant folding
\end{tabular}

\tcblower

\textbf{Official Key:} \textbf{\large B} \hfill \textit{\small [Easy]}
\vspace{2pt}\hrule height 0.3pt\vspace{3pt}
\textbf{Topic:} Compiler Design \\
\textbf{Subtopic:} \textit{Data flow analyses: constant propagation, liveness analysis, common sub expression elimination}
\vspace{2pt}\hrule height 0.3pt\vspace{3pt}
\textbf{Guide / Hint:} \\
{\small Live variable analysis determines whether the value in a variable could be read along some execution path before being redefined. This information is directly used in register allocation and assignment to construct interference graphs and allocate registers efficiently.}

\end{tcolorbox}
\vspace{6pt}\hrule height 0.35pt\vspace{6pt}
\noindent
\begin{tcolorbox}[
    blanker,
    sidebyside,
    sidebyside align=top,
    lefthand width=0.55\textwidth,
    top=4pt, bottom=4pt, left=0pt, right=0pt,
    middle=16pt,
    underlay={
        \draw[dashed, black!80, line width=0.8pt, dash pattern=on 4pt off 3pt]
            ([xshift=0.55\textwidth+8pt]frame.north west) --
            ([xshift=0.55\textwidth+8pt]frame.south west);
    }
]
\textbf{Q.14} \quad {\footnotesize\textsc{[MCQ $\cdot$ 1 Mark(s) $\cdot$ Neg: 0.33]}} \hfill {\tiny\texttt{2025\_CS1\_CS\_Q14}}
\vspace{2pt}\hrule height 0.3pt\vspace{4pt}

Consider a demand paging memory management system with 32-bit logical address, 20-bit physical address, and page size of 2048 bytes. Assuming that the memory is byte addressable, what is the maximum number of entries in the page table?

\par\vspace{6pt}\noindent
\begin{tabular}{@{}p{0.48\linewidth}@{\hspace{0.04\linewidth}}p{0.48\linewidth}@{}}
$\square$\ \textbf{(A)} $2^{21}$ & $\square$\ \textbf{(B)} $2^{20}$ \\
\addlinespace[3pt]
$\square$\ \textbf{(C)} $2^{22}$ & $\square$\ \textbf{(D)} $2^{24}$
\end{tabular}

\tcblower

\textbf{Official Key:} \textbf{\large A} \hfill \textit{\small [Medium]}
\vspace{2pt}\hrule height 0.3pt\vspace{3pt}
\textbf{Topic:} Operating System \\
\textbf{Subtopic:} \textit{Memory management and virtual memory}
\vspace{2pt}\hrule height 0.3pt\vspace{3pt}
\textbf{Guide / Hint:} \\
{\small With a page size of $2048 = 2^{11}$ bytes, the page offset requires 11 bits. The 32-bit logical address is split into page number and page offset, so the page number has $32 - 11 = 21$ bits. Hence, the page table contains at most $2^{21}$ entries.}

\end{tcolorbox}
\vspace{6pt}\hrule height 0.35pt\vspace{6pt}
\noindent
\begin{tcolorbox}[
    blanker,
    sidebyside,
    sidebyside align=top,
    lefthand width=0.55\textwidth,
    top=4pt, bottom=4pt, left=0pt, right=0pt,
    middle=16pt,
    underlay={
        \draw[dashed, black!80, line width=0.8pt, dash pattern=on 4pt off 3pt]
            ([xshift=0.55\textwidth+8pt]frame.north west) --
            ([xshift=0.55\textwidth+8pt]frame.south west);
    }
]
\textbf{Q.15} \quad {\footnotesize\textsc{[MCQ $\cdot$ 1 Mark(s) $\cdot$ Neg: 0.33]}} \hfill {\tiny\texttt{2025\_CS1\_CS\_Q15}}
\vspace{2pt}\hrule height 0.3pt\vspace{4pt}

A schedule of three database transactions $T_1, T_2$, and $T_3$ is shown. $R_i(A)$ and $W_i(A)$ denote read and write of data item $A$ by transaction $T_i$, $i= 1,2,3$. The transaction $T_1$ aborts at the end. Which other transaction(s) will be required to be rolled back?
\[ R\_1(X) \; W\_1(Y) \; R\_2(X) \; R\_2(Y) \; R\_3(Y) \; ABORT(T\_1) \]

\par\vspace{6pt}\noindent
\begin{tabular}{@{}p{0.48\linewidth}@{\hspace{0.04\linewidth}}p{0.48\linewidth}@{}}
$\square$\ \textbf{(A)} Only $T_2$ & $\square$\ \textbf{(B)} Only $T_3$ \\
\addlinespace[3pt]
$\square$\ \textbf{(C)} Both $T_2$ and $T_3$ & $\square$\ \textbf{(D)} Neither $T_2$ nor $T_3$
\end{tabular}

\tcblower

\textbf{Official Key:} \textbf{\large C} \hfill \textit{\small [Medium]}
\vspace{2pt}\hrule height 0.3pt\vspace{3pt}
\textbf{Topic:} Databases \\
\textbf{Subtopic:} \textit{Transactions and concurrency control}
\vspace{2pt}\hrule height 0.3pt\vspace{3pt}
\textbf{Guide / Hint:} \\
{\small Transaction $T_1$ writes to item $Y$ ($W_1(Y)$). Subsequently, both $T_2$ and $T_3$ read uncommitted data item $Y$ ($R_2(Y)$ and $R_3(Y)$) written by $T_1$ before $T_1$ commits. When $T_1$ aborts, a cascading rollback is triggered, requiring both $T_2$ and $T_3$ to be rolled back.}

\end{tcolorbox}
\vspace{6pt}\hrule height 0.35pt\vspace{6pt}
\noindent
\begin{tcolorbox}[
    blanker,
    sidebyside,
    sidebyside align=top,
    lefthand width=0.55\textwidth,
    top=4pt, bottom=4pt, left=0pt, right=0pt,
    middle=16pt,
    underlay={
        \draw[dashed, black!80, line width=0.8pt, dash pattern=on 4pt off 3pt]
            ([xshift=0.55\textwidth+8pt]frame.north west) --
            ([xshift=0.55\textwidth+8pt]frame.south west);
    }
]
\textbf{Q.16} \quad {\footnotesize\textsc{[MCQ $\cdot$ 1 Mark(s) $\cdot$ Neg: 0.33]}} \hfill {\tiny\texttt{2025\_CS1\_CS\_Q16}}
\vspace{2pt}\hrule height 0.3pt\vspace{4pt}

Identify the ONE CORRECT matching between the OSI layers and their corresponding functionalities as shown.
\par\medskip
\begin{tabular}{ll}
\textbf{OSI Layers} & \textbf{Functionalities} \\
(a) Network layer & (I) Packet routing \\
(b) Transport layer & (II) Framing and error handling \\
(c) Datalink layer & (III) Host to host communication
\end{tabular}

\par\vspace{6pt}\noindent
\begin{tabular}{@{}p{0.48\linewidth}@{\hspace{0.04\linewidth}}p{0.48\linewidth}@{}}
$\square$\ \textbf{(A)} (a)-(I), (b)-(II), (c)-(III) & $\square$\ \textbf{(B)} (a)-(I), (b)-(III), (c)-(II) \\
\addlinespace[3pt]
$\square$\ \textbf{(C)} (a)-(II), (b)-(I), (c)-(III) & $\square$\ \textbf{(D)} (a)-(III), (b)-(II), (c)-(I)
\end{tabular}

\tcblower

\textbf{Official Key:} \textbf{\large B} \hfill \textit{\small [Easy]}
\vspace{2pt}\hrule height 0.3pt\vspace{3pt}
\textbf{Topic:} Computer Networks \\
\textbf{Subtopic:} \textit{Concept of layering: OSI and TCP/IP Protocol Stacks}
\vspace{2pt}\hrule height 0.3pt\vspace{3pt}
\textbf{Guide / Hint:} \\
{\small The Network layer is responsible for packet routing across networks ((a)-(I)). The Transport layer provides end-to-end communication services between hosts ((b)-(III)). The Data Link layer handles framing, flow control, and physical error handling ((c)-(II)).}

\end{tcolorbox}
\vspace{6pt}\hrule height 0.35pt\vspace{6pt}
\noindent
\begin{tcolorbox}[
    blanker,
    sidebyside,
    sidebyside align=top,
    lefthand width=0.55\textwidth,
    top=4pt, bottom=4pt, left=0pt, right=0pt,
    middle=16pt,
    underlay={
        \draw[dashed, black!80, line width=0.8pt, dash pattern=on 4pt off 3pt]
            ([xshift=0.55\textwidth+8pt]frame.north west) --
            ([xshift=0.55\textwidth+8pt]frame.south west);
    }
]
\textbf{Q.17} \quad {\footnotesize\textsc{[MCQ $\cdot$ 1 Mark(s) $\cdot$ Neg: 0.33]}} \hfill {\tiny\texttt{2025\_CS1\_CS\_Q17}}
\vspace{2pt}\hrule height 0.3pt\vspace{4pt}

$g(\cdot)$ is a function from $A$ to $B$, $f(\cdot)$ is a function from $B$ to $C$, and their composition defined as $f(g(\cdot))$ is a mapping from $A$ to $C$.
If $f(\cdot)$ and $f(g(\cdot))$ are onto (surjective) functions, which ONE of the following is TRUE about the function $g(\cdot)$?

\par\vspace{6pt}\noindent
\begin{tabular}{@{}p{0.98\linewidth}@{}}
$\square$\ \textbf{(A)} $g(\cdot)$ must be an onto (surjective) function. \\
\addlinespace[3pt]
$\square$\ \textbf{(B)} $g(\cdot)$ must be a one-to-one (injective) function. \\
\addlinespace[3pt]
$\square$\ \textbf{(C)} $g(\cdot)$ must be a bijective function, that is, both one-to-one and onto. \\
\addlinespace[3pt]
$\square$\ \textbf{(D)} $g(\cdot)$ is not required to be a one-to-one or onto function.
\end{tabular}

\tcblower

\textbf{Official Key:} \textbf{\large D} \hfill \textit{\small [Medium]}
\vspace{2pt}\hrule height 0.3pt\vspace{3pt}
\textbf{Topic:} Discrete Mathematics \\
\textbf{Subtopic:} \textit{Sets, relations, functions, partial orders and lattices}
\vspace{2pt}\hrule height 0.3pt\vspace{3pt}
\textbf{Guide / Hint:} \\
{\small If the composition $f \circ g: A \to C$ is surjective, $f$ must be surjective, but $g$ is not required to be either injective or surjective. For example, if $C$ is a singleton $\{c\}$ and $B = \{b_1, b_2\}$, $g$ can map all of $A$ to $b_1$ where $f(b_1) = c$, leaving $b_2$ uncovered (not onto) and having many collisions (not one-to-one).}

\end{tcolorbox}
\vspace{6pt}\hrule height 0.35pt\vspace{6pt}
\noindent
\begin{tcolorbox}[
    blanker,
    sidebyside,
    sidebyside align=top,
    lefthand width=0.55\textwidth,
    top=4pt, bottom=4pt, left=0pt, right=0pt,
    middle=16pt,
    underlay={
        \draw[dashed, black!80, line width=0.8pt, dash pattern=on 4pt off 3pt]
            ([xshift=0.55\textwidth+8pt]frame.north west) --
            ([xshift=0.55\textwidth+8pt]frame.south west);
    }
]
\textbf{Q.18} \quad {\footnotesize\textsc{[MCQ $\cdot$ 1 Mark(s) $\cdot$ Neg: 0.33]}} \hfill {\tiny\texttt{2025\_CS1\_CS\_Q18}}
\vspace{2pt}\hrule height 0.3pt\vspace{4pt}

Let $G$ be any undirected graph with positive edge weights, and $T$ be a minimum spanning tree of $G$. For any two vertices, $u$ and $v$, let $d_1(u, v)$ and $d_2(u, v)$ be the shortest distances between $u$ and $v$ in $G$ and $T$, respectively. Which ONE of the options is CORRECT for all possible $G, T, u$ and $v$?

\par\vspace{6pt}\noindent
\begin{tabular}{@{}p{0.48\linewidth}@{\hspace{0.04\linewidth}}p{0.48\linewidth}@{}}
$\square$\ \textbf{(A)} $d_1(u, v) = d_2(u, v)$ & $\square$\ \textbf{(B)} $d_1(u, v) \le d_2(u, v)$ \\
\addlinespace[3pt]
$\square$\ \textbf{(C)} $d_1(u, v) \ge d_2(u, v)$ & $\square$\ \textbf{(D)} $d_1(u, v) \ne d_2(u, v)$
\end{tabular}

\tcblower

\textbf{Official Key:} \textbf{\large B} \hfill \textit{\small [Easy]}
\vspace{2pt}\hrule height 0.3pt\vspace{3pt}
\textbf{Topic:} Algorithms \\
\textbf{Subtopic:} \textit{Graph traversals, minimum spanning trees, shortest paths}
\vspace{2pt}\hrule height 0.3pt\vspace{3pt}
\textbf{Guide / Hint:} \\
{\small Since the minimum spanning tree $T$ is a subgraph of $G$ ($E(T) \subseteq E(G)$), every path between $u$ and $v$ in $T$ is also a valid path in $G$. Thus, the shortest path distance in $G$ cannot exceed that in $T$, guaranteeing $d_1(u, v) \le d_2(u, v)$.}

\end{tcolorbox}
\vspace{6pt}\hrule height 0.35pt\vspace{6pt}
\noindent
\begin{tcolorbox}[
    blanker,
    sidebyside,
    sidebyside align=top,
    lefthand width=0.55\textwidth,
    top=4pt, bottom=4pt, left=0pt, right=0pt,
    middle=16pt,
    underlay={
        \draw[dashed, black!80, line width=0.8pt, dash pattern=on 4pt off 3pt]
            ([xshift=0.55\textwidth+8pt]frame.north west) --
            ([xshift=0.55\textwidth+8pt]frame.south west);
    }
]
\textbf{Q.19} \quad {\footnotesize\textsc{[MCQ $\cdot$ 1 Mark(s) $\cdot$ Neg: 0.33]}} \hfill {\tiny\texttt{2025\_CS1\_CS\_Q19}}
\vspace{2pt}\hrule height 0.3pt\vspace{4pt}

Consider the following context-free grammar $G$, where $S, A$, and $B$ are the variables (non-terminals), $a$ and $b$ are the terminal symbols, $S$ is the start variable, and the rules of $G$ are described as:
\[ S \to aaB \mid Abb \]
\[ A \to a \mid aA \]
\[ B \to b \mid bB \]
Which ONE of the languages $L(G)$ is accepted by $G$?

\par\vspace{6pt}\noindent
\begin{tabular}{@{}p{0.98\linewidth}@{}}
$\square$\ \textbf{(A)} $L(G) = \{a^2b^n \mid n \ge 1\} \cup \{a^nb^2 \mid n \ge 1\}$ \\
\addlinespace[3pt]
$\square$\ \textbf{(B)} $L(G) = \{a^nb^{2n} \mid n \ge 1\} \cup \{a^{2n}b^n \mid n \ge 1\}$ \\
\addlinespace[3pt]
$\square$\ \textbf{(C)} $L(G) = \{a^nb^n \mid n \ge 1\}$ \\
\addlinespace[3pt]
$\square$\ \textbf{(D)} $L(G) = \{a^{2n}b^{2n} \mid n \ge 1\}$
\end{tabular}

\tcblower

\textbf{Official Key:} \textbf{\large A} \hfill \textit{\small [Easy]}
\vspace{2pt}\hrule height 0.3pt\vspace{3pt}
\textbf{Topic:} Theory of Computation \\
\textbf{Subtopic:} \textit{Context-free grammars and push-down automata}
\vspace{2pt}\hrule height 0.3pt\vspace{3pt}
\textbf{Guide / Hint:} \\
{\small From the rules, $A$ generates $a^+$ (any $a^n$ with $n \ge 1$) and $B$ generates $b^+$ (any $b^n$ with $n \ge 1$). Thus, $aaB$ produces $a^2b^n$ ($n \ge 1$), and $Abb$ produces $a^nb^2$ ($n \ge 1$). Therefore, $L(G) = \{a^2b^n \mid n \ge 1\} \cup \{a^nb^2 \mid n \ge 1\}$.}

\end{tcolorbox}
\vspace{6pt}\hrule height 0.35pt\vspace{6pt}
\noindent
\begin{tcolorbox}[
    blanker,
    sidebyside,
    sidebyside align=top,
    lefthand width=0.55\textwidth,
    top=4pt, bottom=4pt, left=0pt, right=0pt,
    middle=16pt,
    underlay={
        \draw[dashed, black!80, line width=0.8pt, dash pattern=on 4pt off 3pt]
            ([xshift=0.55\textwidth+8pt]frame.north west) --
            ([xshift=0.55\textwidth+8pt]frame.south west);
    }
]
\textbf{Q.20} \quad {\footnotesize\textsc{[MCQ $\cdot$ 1 Mark(s) $\cdot$ Neg: 0.33]}} \hfill {\tiny\texttt{2025\_CS1\_CS\_Q20}}
\vspace{2pt}\hrule height 0.3pt\vspace{4pt}

Consider the following recurrence relation:
\[ T(n) = 2T(n-1) + n2^n \quad \text{for } n > 0, \; T(0) = 1. \]
Which ONE of the following options is CORRECT?

\par\vspace{6pt}\noindent
\begin{tabular}{@{}p{0.48\linewidth}@{\hspace{0.04\linewidth}}p{0.48\linewidth}@{}}
$\square$\ \textbf{(A)} $T(n) = \Theta(n^2 2^n)$ & $\square$\ \textbf{(B)} $T(n) = \Theta(n 2^n)$ \\
\addlinespace[3pt]
$\square$\ \textbf{(C)} $T(n) = \Theta((\log n)^2 2^n)$ & $\square$\ \textbf{(D)} $T(n) = \Theta(4^n)$
\end{tabular}

\tcblower

\textbf{Official Key:} \textbf{\large A} \hfill \textit{\small [Medium]}
\vspace{2pt}\hrule height 0.3pt\vspace{3pt}
\textbf{Topic:} Algorithms \\
\textbf{Subtopic:} \textit{Asymptotic worst case time and space complexity}
\vspace{2pt}\hrule height 0.3pt\vspace{3pt}
\textbf{Guide / Hint:} \\
{\small Dividing both sides by $2^n$ yields $\frac{T(n)}{2^n} = \frac{T(n-1)}{2^{n-1}} + n$. Letting $S(n) = \frac{T(n)}{2^n}$, we have $S(n) = S(n-1) + n = S(0) + \sum_{i=1}^n i = 1 + \frac{n(n+1)}{2} = \Theta(n^2)$. Multiplying back gives $T(n) = \Theta(n^2 2^n)$.}

\end{tcolorbox}
\vspace{6pt}\hrule height 0.35pt\vspace{6pt}
\noindent
\begin{tcolorbox}[
    blanker,
    sidebyside,
    sidebyside align=top,
    lefthand width=0.55\textwidth,
    top=4pt, bottom=4pt, left=0pt, right=0pt,
    middle=16pt,
    underlay={
        \draw[dashed, black!80, line width=0.8pt, dash pattern=on 4pt off 3pt]
            ([xshift=0.55\textwidth+8pt]frame.north west) --
            ([xshift=0.55\textwidth+8pt]frame.south west);
    }
]
\textbf{Q.21} \quad {\footnotesize\textsc{[MSQ $\cdot$ 1 Mark(s) $\cdot$ No Neg]}} \hfill {\tiny\texttt{2025\_CS1\_CS\_Q21}}
\vspace{2pt}\hrule height 0.3pt\vspace{4pt}

Consider a $B^+$ tree with 5 nodes, in which a node can store at most 3 key values. The value 23 is now inserted in the $B^+$ tree. Which of the following option(s) is/are CORRECT?

\par\vspace{6pt}\noindent
\begin{tabular}{@{}p{0.98\linewidth}@{}}
$\square$\ \textbf{(A)} None of the nodes will split. \\
\addlinespace[3pt]
$\square$\ \textbf{(B)} At least one node will split and redistribute. \\
\addlinespace[3pt]
$\square$\ \textbf{(C)} The total number of nodes will remain same. \\
\addlinespace[3pt]
$\square$\ \textbf{(D)} The height of the tree will increase.
\end{tabular}

\tcblower

\textbf{Official Key:} \textbf{\large B;D} \hfill \textit{\small [Medium]}
\vspace{2pt}\hrule height 0.3pt\vspace{3pt}
\textbf{Topic:} Databases \\
\textbf{Subtopic:} \textit{File organization, indexing (e.g., B and B+ trees)}
\vspace{2pt}\hrule height 0.3pt\vspace{3pt}
\textbf{Guide / Hint:} \\
{\small With maximum key capacity of 3 keys per node, the targeted leaf node containing keys $[20, 21, 22]$ is already full. Inserting 23 causes this leaf node to overflow and split, propagating a split upward to the already-saturated root, which creates a new root node and increases tree height.}

\end{tcolorbox}
\vspace{6pt}\hrule height 0.35pt\vspace{6pt}
\noindent
\begin{tcolorbox}[
    blanker,
    sidebyside,
    sidebyside align=top,
    lefthand width=0.55\textwidth,
    top=4pt, bottom=4pt, left=0pt, right=0pt,
    middle=16pt,
    underlay={
        \draw[dashed, black!80, line width=0.8pt, dash pattern=on 4pt off 3pt]
            ([xshift=0.55\textwidth+8pt]frame.north west) --
            ([xshift=0.55\textwidth+8pt]frame.south west);
    }
]
\textbf{Q.22} \quad {\footnotesize\textsc{[MSQ $\cdot$ 1 Mark(s) $\cdot$ No Neg]}} \hfill {\tiny\texttt{2025\_CS1\_CS\_Q22}}
\vspace{2pt}\hrule height 0.3pt\vspace{4pt}

Consider the 3-way handshaking protocol for TCP connection establishment. Let the three packets exchanged during the connection establishment be denoted as P1, P2, and P3, in order. Which of the following option(s) is/are TRUE with respect to TCP header flags that are set in the packets?

\par\vspace{6pt}\noindent
\begin{tabular}{@{}p{0.48\linewidth}@{\hspace{0.04\linewidth}}p{0.48\linewidth}@{}}
$\square$\ \textbf{(A)} P3: SYN = 1, ACK = 1 & $\square$\ \textbf{(B)} P2: SYN = 1, ACK = 1 \\
\addlinespace[3pt]
$\square$\ \textbf{(C)} P2: SYN = 0, ACK = 1 & $\square$\ \textbf{(D)} P1: SYN = 1
\end{tabular}

\tcblower

\textbf{Official Key:} \textbf{\large B;D} \hfill \textit{\small [Easy]}
\vspace{2pt}\hrule height 0.3pt\vspace{3pt}
\textbf{Topic:} Computer Networks \\
\textbf{Subtopic:} \textit{Transport layer: flow control and congestion control, UDP, TCP, sockets}
\vspace{2pt}\hrule height 0.3pt\vspace{3pt}
\textbf{Guide / Hint:} \\
{\small In TCP 3-way handshaking: Packet 1 (Client to Server) sends a SYN segment with flag $[\text{SYN}=1, \text{ACK}=0]$; Packet 2 (Server to Client) responds with $[\text{SYN}=1, \text{ACK}=1]$; Packet 3 (Client to Server) completes setup with $[\text{SYN}=0, \text{ACK}=1]$. Hence, P1 has SYN = 1 and P2 has SYN = 1, ACK = 1.}

\end{tcolorbox}
\vspace{6pt}\hrule height 0.35pt\vspace{6pt}
\noindent
\begin{tcolorbox}[
    blanker,
    sidebyside,
    sidebyside align=top,
    lefthand width=0.55\textwidth,
    top=4pt, bottom=4pt, left=0pt, right=0pt,
    middle=16pt,
    underlay={
        \draw[dashed, black!80, line width=0.8pt, dash pattern=on 4pt off 3pt]
            ([xshift=0.55\textwidth+8pt]frame.north west) --
            ([xshift=0.55\textwidth+8pt]frame.south west);
    }
]
\textbf{Q.23} \quad {\footnotesize\textsc{[MSQ $\cdot$ 1 Mark(s) $\cdot$ No Neg]}} \hfill {\tiny\texttt{2025\_CS1\_CS\_Q23}}
\vspace{2pt}\hrule height 0.3pt\vspace{4pt}

Consider the given system of linear equations for variables $x$ and $y$, where $k$ is a real-valued constant. Which of the following option(s) is/are CORRECT?
\[ x + ky = 1 \]
\[ kx + y = -1 \]

\par\vspace{6pt}\noindent
\begin{tabular}{@{}p{0.98\linewidth}@{}}
$\square$\ \textbf{(A)} There is exactly one value of $k$ for which the above system of equations has no solution. \\
\addlinespace[3pt]
$\square$\ \textbf{(B)} There exist an infinite number of values of $k$ for which the system of equations has no solution. \\
\addlinespace[3pt]
$\square$\ \textbf{(C)} There exists exactly one value of $k$ for which the system of equations has exactly one solution. \\
\addlinespace[3pt]
$\square$\ \textbf{(D)} There exists exactly one value of $k$ for which the system of equations has an infinite number of solutions.
\end{tabular}

\tcblower

\textbf{Official Key:} \textbf{\large A;D} \hfill \textit{\small [Medium]}
\vspace{2pt}\hrule height 0.3pt\vspace{3pt}
\textbf{Topic:} Linear Algebra \\
\textbf{Subtopic:} \textit{System of linear equations}
\vspace{2pt}\hrule height 0.3pt\vspace{3pt}
\textbf{Guide / Hint:} \\
{\small The coefficient determinant is $|A| = 1 - k^2 = 0 \implies k = \pm 1$. For $k = 1$, equations are $x + y = 1$ and $x + y = -1$, which are inconsistent (no solution). For $k = -1$, equations are $x - y = 1$ and $-x + y = -1$, which are identical (infinitely many solutions). For all $k \neq \pm 1$, there is a unique solution.}

\end{tcolorbox}
\vspace{6pt}\hrule height 0.35pt\vspace{6pt}
\noindent
\begin{tcolorbox}[
    blanker,
    sidebyside,
    sidebyside align=top,
    lefthand width=0.55\textwidth,
    top=4pt, bottom=4pt, left=0pt, right=0pt,
    middle=16pt,
    underlay={
        \draw[dashed, black!80, line width=0.8pt, dash pattern=on 4pt off 3pt]
            ([xshift=0.55\textwidth+8pt]frame.north west) --
            ([xshift=0.55\textwidth+8pt]frame.south west);
    }
]
\textbf{Q.24} \quad {\footnotesize\textsc{[MSQ $\cdot$ 1 Mark(s) $\cdot$ No Neg]}} \hfill {\tiny\texttt{2025\_CS1\_CS\_Q24}}
\vspace{2pt}\hrule height 0.3pt\vspace{4pt}

Let $X$ be a 3-variable Boolean function that produces output as `1' when at least two of the input variables are `1'. Which of the following statement(s) is/are CORRECT, where $a, b, c, d, e$ are Boolean variables?

\par\vspace{6pt}\noindent
\begin{tabular}{@{}p{0.98\linewidth}@{}}
$\square$\ \textbf{(A)} $X(a, b, X(c, d, e)) = X(X(a, b, c), d, e)$ \\
\addlinespace[3pt]
$\square$\ \textbf{(B)} $X(a, b, X(a, b, c)) = X(a, b, c)$ \\
\addlinespace[3pt]
$\square$\ \textbf{(C)} $X(a, b, X(a, c, d)) = (X(a, b, a) \text{ AND } X(c, d, c))$ \\
\addlinespace[3pt]
$\square$\ \textbf{(D)} $X(a, b, c) = X(a, X(a, b, c), X(a, c, c))$
\end{tabular}

\tcblower

\textbf{Official Key:} \textbf{\large B;D} \hfill \textit{\small [Medium]}
\vspace{2pt}\hrule height 0.3pt\vspace{3pt}
\textbf{Topic:} Digital Logic \\
\textbf{Subtopic:} \textit{Boolean algebra}
\vspace{2pt}\hrule height 0.3pt\vspace{3pt}
\textbf{Guide / Hint:} \\
{\small The majority function is $X(a, b, c) = ab + bc + ca$. In (B), $X(a, b, M) = ab + bM + aM = ab + (a+b)M = ab + (a+b)(ab+bc+ca) = ab + bc + ca = X(a, b, c)$, so (B) holds. In (D), $X(a, c, c) = c$, giving $X(a, M, c) = aM + Mc + ca = (a+c)M + ac = (a+c)(ab+bc+ca) + ac = ab + bc + ca = X(a, b, c)$, so (D) holds.}

\end{tcolorbox}
\vspace{6pt}\hrule height 0.35pt\vspace{6pt}
\noindent
\begin{tcolorbox}[
    blanker,
    sidebyside,
    sidebyside align=top,
    lefthand width=0.55\textwidth,
    top=4pt, bottom=4pt, left=0pt, right=0pt,
    middle=16pt,
    underlay={
        \draw[dashed, black!80, line width=0.8pt, dash pattern=on 4pt off 3pt]
            ([xshift=0.55\textwidth+8pt]frame.north west) --
            ([xshift=0.55\textwidth+8pt]frame.south west);
    }
]
\textbf{Q.25} \quad {\footnotesize\textsc{[MSQ $\cdot$ 1 Mark(s) $\cdot$ No Neg]}} \hfill {\tiny\texttt{2025\_CS1\_CS\_Q25}}
\vspace{2pt}\hrule height 0.3pt\vspace{4pt}

The number $-6$ can be represented as $1010_2$ in 4-bit 2's complement representation. Which of the following is/are CORRECT 2's complement representation(s) of $-6$?

\par\vspace{6pt}\noindent
\begin{tabular}{@{}p{0.48\linewidth}@{\hspace{0.04\linewidth}}p{0.48\linewidth}@{}}
$\square$\ \textbf{(A)} 1000 1010 in 8-bits & $\square$\ \textbf{(B)} 1111 1010 in 8-bits \\
\addlinespace[3pt]
$\square$\ \textbf{(C)} 1000 0000 0000 1010 in 16-bits & $\square$\ \textbf{(D)} 1111 1111 1111 1010 in 16-bits
\end{tabular}

\tcblower

\textbf{Official Key:} \textbf{\large B;D} \hfill \textit{\small [Easy]}
\vspace{2pt}\hrule height 0.3pt\vspace{3pt}
\textbf{Topic:} Digital Logic \\
\textbf{Subtopic:} \textit{Number representations and computer arithmetic (fixed and floating point)}
\vspace{2pt}\hrule height 0.3pt\vspace{3pt}
\textbf{Guide / Hint:} \\
{\small Sign extension in two's complement preserves the numerical value by replicating the most significant bit (sign bit). Since the 4-bit representation of $-6$ is $1010$ (sign bit 1), extending to 8 bits gives $1111\,1010$ and extending to 16 bits gives $1111\,1111\,1111\,1010$.}

\end{tcolorbox}
\vspace{6pt}\hrule height 0.35pt\vspace{6pt}
\noindent
\begin{tcolorbox}[
    blanker,
    sidebyside,
    sidebyside align=top,
    lefthand width=0.55\textwidth,
    top=4pt, bottom=4pt, left=0pt, right=0pt,
    middle=16pt,
    underlay={
        \draw[dashed, black!80, line width=0.8pt, dash pattern=on 4pt off 3pt]
            ([xshift=0.55\textwidth+8pt]frame.north west) --
            ([xshift=0.55\textwidth+8pt]frame.south west);
    }
]
\textbf{Q.26} \quad {\footnotesize\textsc{[MSQ $\cdot$ 1 Mark(s) $\cdot$ No Neg]}} \hfill {\tiny\texttt{2025\_CS1\_CS\_Q26}}
\vspace{2pt}\hrule height 0.3pt\vspace{4pt}

Which of the following statement(s) is/are TRUE for any binary search tree (BST) having $n$ distinct integers?

\par\vspace{6pt}\noindent
\begin{tabular}{@{}p{0.98\linewidth}@{}}
$\square$\ \textbf{(A)} The maximum length of a path from the root node to any other node is $(n-1)$. \\
\addlinespace[3pt]
$\square$\ \textbf{(B)} An inorder traversal will always produce a sorted sequence of elements. \\
\addlinespace[3pt]
$\square$\ \textbf{(C)} Finding an element takes $O(\log_2 n)$ time in the worst case. \\
\addlinespace[3pt]
$\square$\ \textbf{(D)} Every BST is also a Min-Heap.
\end{tabular}

\tcblower

\textbf{Official Key:} \textbf{\large A;B} \hfill \textit{\small [Easy]}
\vspace{2pt}\hrule height 0.3pt\vspace{3pt}
\textbf{Topic:} Programming and Data Structures \\
\textbf{Subtopic:} \textit{Trees, binary search trees, binary heaps}
\vspace{2pt}\hrule height 0.3pt\vspace{3pt}
\textbf{Guide / Hint:} \\
{\small In a degenerate (skewed) binary search tree, the path length from root to the farthest leaf is $n - 1$ edges (A is true). Inorder traversal of any valid BST visits nodes in strictly increasing order (B is true). Worst-case search in an unbalanced BST takes $O(n)$ time, and BST property differs fundamentally from heap order.}

\end{tcolorbox}
\vspace{6pt}\hrule height 0.35pt\vspace{6pt}
\noindent
\begin{tcolorbox}[
    blanker,
    sidebyside,
    sidebyside align=top,
    lefthand width=0.55\textwidth,
    top=4pt, bottom=4pt, left=0pt, right=0pt,
    middle=16pt,
    underlay={
        \draw[dashed, black!80, line width=0.8pt, dash pattern=on 4pt off 3pt]
            ([xshift=0.55\textwidth+8pt]frame.north west) --
            ([xshift=0.55\textwidth+8pt]frame.south west);
    }
]
\textbf{Q.27} \quad {\footnotesize\textsc{[MSQ $\cdot$ 1 Mark(s) $\cdot$ No Neg]}} \hfill {\tiny\texttt{2025\_CS1\_CS\_Q27}}
\vspace{2pt}\hrule height 0.3pt\vspace{4pt}

A partial data path of a processor is given, where RA, RB, and RZ are 32-bit registers. Multiplexer Mux\_A selects between RA and a 32-bit immediate value, and Mux\_B selects between RB and a 32-bit immediate value. The output of the ALU is stored into RZ. Which option(s) is/are CORRECT related to arithmetic operations using the data path as shown?

\par\vspace{6pt}\noindent
\begin{tabular}{@{}p{0.98\linewidth}@{}}
$\square$\ \textbf{(A)} The data path can implement arithmetic operations involving two registers. \\
\addlinespace[3pt]
$\square$\ \textbf{(B)} The data path can implement arithmetic operations involving one register and one immediate value. \\
\addlinespace[3pt]
$\square$\ \textbf{(C)} The data path can implement arithmetic operations involving two immediate values. \\
\addlinespace[3pt]
$\square$\ \textbf{(D)} The data path can only implement arithmetic operations involving one register and one immediate value.
\end{tabular}

\tcblower

\textbf{Official Key:} \textbf{\large A;B;C} \hfill \textit{\small [Medium]}
\vspace{2pt}\hrule height 0.3pt\vspace{3pt}
\textbf{Topic:} Computer Organization and Architecture \\
\textbf{Subtopic:} \textit{ALU, data-path and control unit}
\vspace{2pt}\hrule height 0.3pt\vspace{3pt}
\textbf{Guide / Hint:} \\
{\small Since Mux\_A can select either register RA or an immediate value, and Mux\_B can independently select either register RB or an immediate value, the ALU inputs can be configured as (RA, RB), (RA, immediate), (immediate, RB), or (immediate, immediate). Thus A, B, and C are all valid.}

\end{tcolorbox}
\vspace{6pt}\hrule height 0.35pt\vspace{6pt}
\noindent
\begin{tcolorbox}[
    blanker,
    sidebyside,
    sidebyside align=top,
    lefthand width=0.55\textwidth,
    top=4pt, bottom=4pt, left=0pt, right=0pt,
    middle=16pt,
    underlay={
        \draw[dashed, black!80, line width=0.8pt, dash pattern=on 4pt off 3pt]
            ([xshift=0.55\textwidth+8pt]frame.north west) --
            ([xshift=0.55\textwidth+8pt]frame.south west);
    }
]
\textbf{Q.28} \quad {\footnotesize\textsc{[MSQ $\cdot$ 1 Mark(s) $\cdot$ No Neg]}} \hfill {\tiny\texttt{2025\_CS1\_CS\_Q28}}
\vspace{2pt}\hrule height 0.3pt\vspace{4pt}

A regular language $L$ is accepted by a non-deterministic finite automaton (NFA) with $n$ states. Which of the following statement(s) is/are FALSE?

\par\vspace{6pt}\noindent
\begin{tabular}{@{}p{0.98\linewidth}@{}}
$\square$\ \textbf{(A)} $L$ may have an accepting NFA with $< n$ states. \\
\addlinespace[3pt]
$\square$\ \textbf{(B)} $L$ may have an accepting DFA with $< n$ states. \\
\addlinespace[3pt]
$\square$\ \textbf{(C)} There exists a DFA with $\le 2^n$ states that accepts $L$. \\
\addlinespace[3pt]
$\square$\ \textbf{(D)} Every DFA that accepts $L$ has $> 2^n$ states.
\end{tabular}

\tcblower

\textbf{Official Key:} \textbf{\large D} \hfill \textit{\small [Easy]}
\vspace{2pt}\hrule height 0.3pt\vspace{3pt}
\textbf{Topic:} Theory of Computation \\
\textbf{Subtopic:} \textit{Regular expressions and finite automata}
\vspace{2pt}\hrule height 0.3pt\vspace{3pt}
\textbf{Guide / Hint:} \\
{\small By the standard subset construction algorithm, every $n$-state NFA can be converted into an equivalent deterministic finite automaton (DFA) having at most $2^n$ states. Statement D asserts that every DFA must have strictly more than $2^n$ states, which is directly false.}

\end{tcolorbox}
\vspace{6pt}\hrule height 0.35pt\vspace{6pt}
\noindent
\begin{tcolorbox}[
    blanker,
    sidebyside,
    sidebyside align=top,
    lefthand width=0.55\textwidth,
    top=4pt, bottom=4pt, left=0pt, right=0pt,
    middle=16pt,
    underlay={
        \draw[dashed, black!80, line width=0.8pt, dash pattern=on 4pt off 3pt]
            ([xshift=0.55\textwidth+8pt]frame.north west) --
            ([xshift=0.55\textwidth+8pt]frame.south west);
    }
]
\textbf{Q.29} \quad {\footnotesize\textsc{[NAT $\cdot$ 1 Mark(s) $\cdot$ No Neg]}} \hfill {\tiny\texttt{2025\_CS1\_CS\_Q29}}
\vspace{2pt}\hrule height 0.3pt\vspace{4pt}

Suppose in a multiprogramming environment, the following C program segment is executed. A process goes into I/O queue whenever an I/O related operation is performed. Assume that there will always be a context switch whenever a process requests for an I/O, and also whenever the process returns from an I/O. The number of times the process will enter the ready queue during its lifetime (not counting the time the process enters the ready queue when it is run initially) is \_\_\_\_\_\_\. (Answer in integer)
\par\medskip
\begin{verbatim}
int main()
{
    int x=0, i=0;
    scanf("%d", &x);
    for(i=0; i<20; i++)
    {
        x = x + 20;
        printf("%d\n", x);
    }
    return 0;
}
\end{verbatim}

\par\vspace{6pt}\noindent\textbf{Your Answer:} \fbox{\rule{0pt}{1.3ex}\hspace{3cm}}

\tcblower

\textbf{Official Key:} \textbf{\large 21 to 21} \hfill \textit{\small [Medium]}
\vspace{2pt}\hrule height 0.3pt\vspace{3pt}
\textbf{Topic:} Operating System \\
\textbf{Subtopic:} \textit{CPU and I/O scheduling}
\vspace{2pt}\hrule height 0.3pt\vspace{3pt}
\textbf{Guide / Hint:} \\
{\small There is 1 initial `scanf` and 20 iterations of `printf`, totaling $1 + 20 = 21$ I/O operations. Each time an I/O operation completes, the process transitions from the I/O (waiting) queue back into the ready queue. Excluding the initial program launch into the ready queue, the process enters the ready queue exactly 21 times.}

\end{tcolorbox}
\vspace{6pt}\hrule height 0.35pt\vspace{6pt}
\noindent
\begin{tcolorbox}[
    blanker,
    sidebyside,
    sidebyside align=top,
    lefthand width=0.55\textwidth,
    top=4pt, bottom=4pt, left=0pt, right=0pt,
    middle=16pt,
    underlay={
        \draw[dashed, black!80, line width=0.8pt, dash pattern=on 4pt off 3pt]
            ([xshift=0.55\textwidth+8pt]frame.north west) --
            ([xshift=0.55\textwidth+8pt]frame.south west);
    }
]
\textbf{Q.30} \quad {\footnotesize\textsc{[NAT $\cdot$ 1 Mark(s) $\cdot$ No Neg]}} \hfill {\tiny\texttt{2025\_CS1\_CS\_Q30}}
\vspace{2pt}\hrule height 0.3pt\vspace{4pt}

Let $S$ be the set of all ternary strings defined over the alphabet $\{a, b, c\}$. Consider all strings in $S$ that contain at least one occurrence of two consecutive symbols, that is, ``aa'', ``bb'' or ``cc''. The number of such strings of length 5 that are possible is \_\_\_\_\_\_\. (Answer in integer)

\par\vspace{6pt}\noindent\textbf{Your Answer:} \fbox{\rule{0pt}{1.3ex}\hspace{3cm}}

\tcblower

\textbf{Official Key:} \textbf{\large 195 to 195} \hfill \textit{\small [Medium]}
\vspace{2pt}\hrule height 0.3pt\vspace{3pt}
\textbf{Topic:} Discrete Mathematics \\
\textbf{Subtopic:} \textit{Combinatorics: counting, recurrence relations, generating functions}
\vspace{2pt}\hrule height 0.3pt\vspace{3pt}
\textbf{Guide / Hint:} \\
{\small The total number of ternary strings of length 5 is $3^5 = 243$. A string has no consecutive identical symbols if the first character has 3 choices and each subsequent character has 2 choices (different from the preceding character), giving $3 \times 2^4 = 48$ such strings. By complementary counting, the number of strings with at least one consecutive pair is $243 - 48 = 195$.}

\end{tcolorbox}
\vspace{6pt}\hrule height 0.35pt\vspace{6pt}
\noindent
\begin{tcolorbox}[
    blanker,
    sidebyside,
    sidebyside align=top,
    lefthand width=0.55\textwidth,
    top=4pt, bottom=4pt, left=0pt, right=0pt,
    middle=16pt,
    underlay={
        \draw[dashed, black!80, line width=0.8pt, dash pattern=on 4pt off 3pt]
            ([xshift=0.55\textwidth+8pt]frame.north west) --
            ([xshift=0.55\textwidth+8pt]frame.south west);
    }
]
\textbf{Q.31} \quad {\footnotesize\textsc{[NAT $\cdot$ 1 Mark(s) $\cdot$ No Neg]}} \hfill {\tiny\texttt{2025\_CS1\_CS\_Q31}}
\vspace{2pt}\hrule height 0.3pt\vspace{4pt}

Consider the given function $f(x)$:
\[ f(x) = \begin{cases} ax + b & \text{for } x < 1 \\ x^3 + x^2 + 1 & \text{for } x \ge 1 \end{cases} \]
If the function is differentiable everywhere, the value of $b$ must be \_\_\_\_\_\_\. (rounded off to one decimal place)

\par\vspace{6pt}\noindent\textbf{Your Answer:} \fbox{\rule{0pt}{1.3ex}\hspace{3cm}}

\tcblower

\textbf{Official Key:} \textbf{\large -2.1 to -1.9} \hfill \textit{\small [Easy]}
\vspace{2pt}\hrule height 0.3pt\vspace{3pt}
\textbf{Topic:} Calculus \\
\textbf{Subtopic:} \textit{Limits, continuity and differentiability}
\vspace{2pt}\hrule height 0.3pt\vspace{3pt}
\textbf{Guide / Hint:} \\
{\small For differentiability at $x = 1$, left and right derivatives must be equal: $\lim_{x \to 1^-} f'(x) = a$ and $\lim_{x \to 1^+} f'(x) = 3(1)^2 + 2(1) = 5$, so $a = 5$. Continuity at $x = 1$ requires $a(1) + b = 1^3 + 1^2 + 1 = 3$. Substituting $a = 5$ gives $5 + b = 3 \implies b = -2.0$.}

\end{tcolorbox}
\vspace{6pt}\hrule height 0.35pt\vspace{6pt}
\noindent
\begin{tcolorbox}[
    blanker,
    sidebyside,
    sidebyside align=top,
    lefthand width=0.55\textwidth,
    top=4pt, bottom=4pt, left=0pt, right=0pt,
    middle=16pt,
    underlay={
        \draw[dashed, black!80, line width=0.8pt, dash pattern=on 4pt off 3pt]
            ([xshift=0.55\textwidth+8pt]frame.north west) --
            ([xshift=0.55\textwidth+8pt]frame.south west);
    }
]
\textbf{Q.32} \quad {\footnotesize\textsc{[NAT $\cdot$ 1 Mark(s) $\cdot$ No Neg]}} \hfill {\tiny\texttt{2025\_CS1\_CS\_Q32}}
\vspace{2pt}\hrule height 0.3pt\vspace{4pt}

A box contains 5 coins: 4 regular coins and 1 fake coin. When a regular coin is tossed, the probability $P(\text{head}) = 0.5$ and for a fake coin, $P(\text{head}) = 1$. You pick a coin at random and toss it twice, and get two heads. The probability that the coin you have chosen is the fake coin is \_\_\_\_\_\_\. (rounded off to two decimal places)

\par\vspace{6pt}\noindent\textbf{Your Answer:} \fbox{\rule{0pt}{1.3ex}\hspace{3cm}}

\tcblower

\textbf{Official Key:} \textbf{\large 0.49 to 0.51} \hfill \textit{\small [Medium]}
\vspace{2pt}\hrule height 0.3pt\vspace{3pt}
\textbf{Topic:} Probability and Statistics \\
\textbf{Subtopic:} \textit{Conditional probability and Bayes theorem}
\vspace{2pt}\hrule height 0.3pt\vspace{3pt}
\textbf{Guide / Hint:} \\
{\small Let $F$ be the event of picking the fake coin ($P(F) = 1/5$) and $R$ picking a regular coin ($P(R) = 4/5$). The likelihood of two heads is $P(2H \mid F) = 1^2 = 1$ and $P(2H \mid R) = (0.5)^2 = 0.25$. By Bayes' theorem, $P(F \mid 2H) = \frac{(1/5)(1)}{(1/5)(1) + (4/5)(0.25)} = \frac{0.2}{0.2 + 0.2} = 0.50$.}

\end{tcolorbox}
\vspace{6pt}\hrule height 0.35pt\vspace{6pt}
\noindent
\begin{tcolorbox}[
    blanker,
    sidebyside,
    sidebyside align=top,
    lefthand width=0.55\textwidth,
    top=4pt, bottom=4pt, left=0pt, right=0pt,
    middle=16pt,
    underlay={
        \draw[dashed, black!80, line width=0.8pt, dash pattern=on 4pt off 3pt]
            ([xshift=0.55\textwidth+8pt]frame.north west) --
            ([xshift=0.55\textwidth+8pt]frame.south west);
    }
]
\textbf{Q.33} \quad {\footnotesize\textsc{[NAT $\cdot$ 1 Mark(s) $\cdot$ No Neg]}} \hfill {\tiny\texttt{2025\_CS1\_CS\_Q33}}
\vspace{2pt}\hrule height 0.3pt\vspace{4pt}

The pseudocode of a function fun() is given below:
\begin{verbatim}
fun(int A[0,...,n-1]){
    for i=0 to n-2
        for j=0 to n-i-2
            if (A[j] > A[j+1])
                then swap A[j] and A[j+1]
}
\end{verbatim}
Let $A[0, \ldots, 29]$ be an array storing 30 distinct integers in descending order. The number of swap operations that will be performed, if the function fun() is called with $A[0, \ldots, 29]$ as argument, is \_\_\_\_\_\_\. (Answer in integer)

\par\vspace{6pt}\noindent\textbf{Your Answer:} \fbox{\rule{0pt}{1.3ex}\hspace{3cm}}

\tcblower

\textbf{Official Key:} \textbf{\large 435 to 435} \hfill \textit{\small [Easy]}
\vspace{2pt}\hrule height 0.3pt\vspace{3pt}
\textbf{Topic:} Algorithms \\
\textbf{Subtopic:} \textit{Searching, sorting, hashing}
\vspace{2pt}\hrule height 0.3pt\vspace{3pt}
\textbf{Guide / Hint:} \\
{\small The function implements standard Bubble Sort. When sorting an array of $n = 30$ elements in strictly reverse order, every comparison triggers a swap operation. The total number of swaps is $\sum_{i=0}^{n-2} (n - 1 - i) = \frac{n(n-1)}{2} = \frac{30 \times 29}{2} = 435$.}

\end{tcolorbox}
\vspace{6pt}\hrule height 0.35pt\vspace{6pt}
\noindent
\begin{tcolorbox}[
    blanker,
    sidebyside,
    sidebyside align=top,
    lefthand width=0.55\textwidth,
    top=4pt, bottom=4pt, left=0pt, right=0pt,
    middle=16pt,
    underlay={
        \draw[dashed, black!80, line width=0.8pt, dash pattern=on 4pt off 3pt]
            ([xshift=0.55\textwidth+8pt]frame.north west) --
            ([xshift=0.55\textwidth+8pt]frame.south west);
    }
]
\textbf{Q.34} \quad {\footnotesize\textsc{[NAT $\cdot$ 1 Mark(s) $\cdot$ No Neg]}} \hfill {\tiny\texttt{2025\_CS1\_CS\_Q34}}
\vspace{2pt}\hrule height 0.3pt\vspace{4pt}

Consider the following C program:
\par\medskip
\begin{verbatim}
#include <stdio.h>
void foo(int *p, int x){
    *p = x;
}
int main(){
    int *z;
    int a = 20, b = 25;
    z = &a;
    foo(z, b);
    printf("%d", a);
    return 0;
}
\end{verbatim}
The output of the given C program is \_\_\_\_\_\_\. (Answer in integer)

\par\vspace{6pt}\noindent\textbf{Your Answer:} \fbox{\rule{0pt}{1.3ex}\hspace{3cm}}

\tcblower

\textbf{Official Key:} \textbf{\large 25 to 25} \hfill \textit{\small [Easy]}
\vspace{2pt}\hrule height 0.3pt\vspace{3pt}
\textbf{Topic:} Programming and Data Structures \\
\textbf{Subtopic:} \textit{Programming in C}
\vspace{2pt}\hrule height 0.3pt\vspace{3pt}
\textbf{Guide / Hint:} \\
{\small In `main()`, pointer `z` is assigned the memory address of variable `a` (`\&a`). Function `foo(z, b)` dereferences `p` (which points to `a`) and assigns `x = b = 25` to `*p`. Thus, `a` is updated to 25, which `printf` prints.}

\end{tcolorbox}
\vspace{6pt}\hrule height 0.35pt\vspace{6pt}
\noindent
\begin{tcolorbox}[
    blanker,
    sidebyside,
    sidebyside align=top,
    lefthand width=0.55\textwidth,
    top=4pt, bottom=4pt, left=0pt, right=0pt,
    middle=16pt,
    underlay={
        \draw[dashed, black!80, line width=0.8pt, dash pattern=on 4pt off 3pt]
            ([xshift=0.55\textwidth+8pt]frame.north west) --
            ([xshift=0.55\textwidth+8pt]frame.south west);
    }
]
\textbf{Q.35} \quad {\footnotesize\textsc{[NAT $\cdot$ 1 Mark(s) $\cdot$ No Neg]}} \hfill {\tiny\texttt{2025\_CS1\_CS\_Q35}}
\vspace{2pt}\hrule height 0.3pt\vspace{4pt}

The height of any rooted tree is defined as the maximum number of edges in the path from the root node to any leaf node. Suppose a Min-Heap $T$ stores 32 keys. The height of $T$ is \_\_\_\_\_\_\. (Answer in integer)

\par\vspace{6pt}\noindent\textbf{Your Answer:} \fbox{\rule{0pt}{1.3ex}\hspace{3cm}}

\tcblower

\textbf{Official Key:} \textbf{\large 5 to 5} \hfill \textit{\small [Easy]}
\vspace{2pt}\hrule height 0.3pt\vspace{3pt}
\textbf{Topic:} Programming and Data Structures \\
\textbf{Subtopic:} \textit{Trees, binary search trees, binary heaps}
\vspace{2pt}\hrule height 0.3pt\vspace{3pt}
\textbf{Guide / Hint:} \\
{\small A binary heap storing $n$ keys is an almost complete binary tree whose height in terms of edges is $\lfloor \log_2 n \rfloor$. For $n = 32$, the height is $\lfloor \log_2 32 \rfloor = 5$.}

\end{tcolorbox}
\vspace{6pt}\hrule height 0.35pt\vspace{6pt}
\noindent
\begin{tcolorbox}[
    blanker,
    sidebyside,
    sidebyside align=top,
    lefthand width=0.55\textwidth,
    top=4pt, bottom=4pt, left=0pt, right=0pt,
    middle=16pt,
    underlay={
        \draw[dashed, black!80, line width=0.8pt, dash pattern=on 4pt off 3pt]
            ([xshift=0.55\textwidth+8pt]frame.north west) --
            ([xshift=0.55\textwidth+8pt]frame.south west);
    }
]
\textbf{Q.36} \quad {\footnotesize\textsc{[MCQ $\cdot$ 2 Mark(s) $\cdot$ Neg: 0.66]}} \hfill {\tiny\texttt{2025\_CS1\_CS\_Q36}}
\vspace{2pt}\hrule height 0.3pt\vspace{4pt}

Consider a memory system with 1M bytes of main memory and 16K bytes of cache memory. Assume that the processor generates 20-bit memory address, and the cache block size is 16 bytes. If the cache uses direct mapping, how many bits will be required to store all the tag values? [Assume memory is byte addressable, $1\text{K}=2^{10}, 1\text{M}=2^{20}$.]

\par\vspace{6pt}\noindent
\begin{tabular}{@{}p{0.48\linewidth}@{\hspace{0.04\linewidth}}p{0.48\linewidth}@{}}
$\square$\ \textbf{(A)} $6 \times 2^{10}$ & $\square$\ \textbf{(B)} $8 \times 2^{10}$ \\
\addlinespace[3pt]
$\square$\ \textbf{(C)} $2^{12}$ & $\square$\ \textbf{(D)} $2^{14}$
\end{tabular}

\tcblower

\textbf{Official Key:} \textbf{\large A} \hfill \textit{\small [Medium]}
\vspace{2pt}\hrule height 0.3pt\vspace{3pt}
\textbf{Topic:} Computer Organization and Architecture \\
\textbf{Subtopic:} \textit{Memory hierarchy: cache, main memory and secondary storage}
\vspace{2pt}\hrule height 0.3pt\vspace{3pt}
\textbf{Guide / Hint:} \\
{\small With block size $16 = 2^4$ bytes, block offset is 4 bits. Cache size is $16\text{ KB} = 2^{14}$ bytes, so number of cache lines is $2^{14}/2^4 = 2^{10} = 1024$, requiring 10 line index bits. The tag bits per line are $20 - (10 + 4) = 6$ bits. For all $2^{10}$ lines, total tag bits required are $6 \times 2^{10}$.}

\end{tcolorbox}
\vspace{6pt}\hrule height 0.35pt\vspace{6pt}
\noindent
\begin{tcolorbox}[
    blanker,
    sidebyside,
    sidebyside align=top,
    lefthand width=0.55\textwidth,
    top=4pt, bottom=4pt, left=0pt, right=0pt,
    middle=16pt,
    underlay={
        \draw[dashed, black!80, line width=0.8pt, dash pattern=on 4pt off 3pt]
            ([xshift=0.55\textwidth+8pt]frame.north west) --
            ([xshift=0.55\textwidth+8pt]frame.south west);
    }
]
\textbf{Q.37} \quad {\footnotesize\textsc{[MCQ $\cdot$ 2 Mark(s) $\cdot$ Neg: 0.66]}} \hfill {\tiny\texttt{2025\_CS1\_CS\_Q37}}
\vspace{2pt}\hrule height 0.3pt\vspace{4pt}

A processor has 64 general-purpose registers and 50 distinct instruction types. An instruction is encoded in 32-bits. What is the maximum number of bits that can be used to store the immediate operand for the given instruction?
\begin{verbatim}
ADD R1, #25 // R1 = R1 + 25
\end{verbatim}

\par\vspace{6pt}\noindent
\begin{tabular}{@{}p{0.48\linewidth}@{\hspace{0.04\linewidth}}p{0.48\linewidth}@{}}
$\square$\ \textbf{(A)} 16 & $\square$\ \textbf{(B)} 20 \\
\addlinespace[3pt]
$\square$\ \textbf{(C)} 22 & $\square$\ \textbf{(D)} 24
\end{tabular}

\tcblower

\textbf{Official Key:} \textbf{\large B} \hfill \textit{\small [Medium]}
\vspace{2pt}\hrule height 0.3pt\vspace{3pt}
\textbf{Topic:} Computer Organization and Architecture \\
\textbf{Subtopic:} \textit{Machine instructions and addressing modes}
\vspace{2pt}\hrule height 0.3pt\vspace{3pt}
\textbf{Guide / Hint:} \\
{\small To represent 50 distinct instruction opcodes, $\lceil \log_2 50 \rceil = 6$ opcode bits are required. Specifying one destination register among 64 general-purpose registers requires $\lceil \log_2 64 \rceil = 6$ bits. In a 32-bit instruction word, the remaining bits for the immediate field are $32 - (6 + 6) = 20$ bits.}

\end{tcolorbox}
\vspace{6pt}\hrule height 0.35pt\vspace{6pt}
\noindent
\begin{tcolorbox}[
    blanker,
    sidebyside,
    sidebyside align=top,
    lefthand width=0.55\textwidth,
    top=4pt, bottom=4pt, left=0pt, right=0pt,
    middle=16pt,
    underlay={
        \draw[dashed, black!80, line width=0.8pt, dash pattern=on 4pt off 3pt]
            ([xshift=0.55\textwidth+8pt]frame.north west) --
            ([xshift=0.55\textwidth+8pt]frame.south west);
    }
]
\textbf{Q.38} \quad {\footnotesize\textsc{[MCQ $\cdot$ 2 Mark(s) $\cdot$ Neg: 0.66]}} \hfill {\tiny\texttt{2025\_CS1\_CS\_Q38}}
\vspace{2pt}\hrule height 0.3pt\vspace{4pt}

A computer has two processors, $M_1$ and $M_2$. Four processes $P_1, P_2, P_3, P_4$ with CPU bursts of 20, 16, 25, and 10 milliseconds, respectively, arrive at the same time and these are the only processes in the system. The scheduler uses non-preemptive priority scheduling, with priorities decided as follows:
\begin{itemize}
\item $M_1$ uses priority of execution for the processes as, $P_1 > P_3 > P_2 > P_4$, i.e., $P_1$ and $P_4$ have highest and lowest priorities, respectively.
\item $M_2$ uses priority of execution for the processes as, $P_2 > P_3 > P_4 > P_1$, i.e., $P_2$ and $P_1$ have highest and lowest priorities, respectively.
\end{itemize}
A process $P_i$ is scheduled to a processor $M_k$, if the processor is free and no other process $P_j$ is waiting with higher priority. At any given point of time, a process can be allocated to any one of the free processors without violating the execution priority rules. Ignore the context switch time. What will be the average waiting time of the processes in milliseconds?

\par\vspace{6pt}\noindent
\begin{tabular}{@{}p{0.48\linewidth}@{\hspace{0.04\linewidth}}p{0.48\linewidth}@{}}
$\square$\ \textbf{(A)} 9.00 & $\square$\ \textbf{(B)} 8.75 \\
\addlinespace[3pt]
$\square$\ \textbf{(C)} 6.50 & $\square$\ \textbf{(D)} 7.50
\end{tabular}

\tcblower

\textbf{Official Key:} \textbf{\large A} \hfill \textit{\small [Hard]}
\vspace{2pt}\hrule height 0.3pt\vspace{3pt}
\textbf{Topic:} Operating System \\
\textbf{Subtopic:} \textit{CPU and I/O scheduling}
\vspace{2pt}\hrule height 0.3pt\vspace{3pt}
\textbf{Guide / Hint:} \\
{\small At $t = 0$, $M_1$ takes highest priority process $P_1$ (burst 20) and $M_2$ takes $P_2$ (burst 16). At $t = 16$, $M_2$ frees up; among remaining processes $\{P_3, P_4\}$, $P_3$ has higher priority on $M_2$, so $P_3$ runs on $M_2$ ($t = 16$ to 41). At $t = 20$, $M_1$ frees up and executes $P_4$ ($t = 20$ to 30). Waiting times are: $P_1 = 0$, $P_2 = 0$, $P_3 = 16$, $P_4 = 20$. Average waiting time is $(0 + 0 + 16 + 20) / 4 = 36 / 4 = 9.00$ ms.}

\end{tcolorbox}
\vspace{6pt}\hrule height 0.35pt\vspace{6pt}
\noindent
\begin{tcolorbox}[
    blanker,
    sidebyside,
    sidebyside align=top,
    lefthand width=0.55\textwidth,
    top=4pt, bottom=4pt, left=0pt, right=0pt,
    middle=16pt,
    underlay={
        \draw[dashed, black!80, line width=0.8pt, dash pattern=on 4pt off 3pt]
            ([xshift=0.55\textwidth+8pt]frame.north west) --
            ([xshift=0.55\textwidth+8pt]frame.south west);
    }
]
\textbf{Q.39} \quad {\footnotesize\textsc{[MCQ $\cdot$ 2 Mark(s) $\cdot$ Neg: 0.66]}} \hfill {\tiny\texttt{2025\_CS1\_CS\_Q39}}
\vspace{2pt}\hrule height 0.3pt\vspace{4pt}

Consider two relations describing $teams$ and $players$ in a sports league:
\begin{itemize}
\item $teams(tid, tname)$: $tid, tname$ are team-id and team-name, respectively
\item $players(pid, pname, tid)$: $pid, pname$, and $tid$ denote player-id, player-name and the team-id of the player, respectively
\end{itemize}
Which ONE of the following tuple relational calculus queries returns the name of the players who play for the team having $tname$ as 'MI'?

\par\vspace{6pt}\noindent
\begin{tabular}{@{}p{0.98\linewidth}@{}}
$\square$\ \textbf{(A)} $\{ p.pname \mid p \in players \land \exists t (t \in teams \land p.tid = t.tid \land t.tname = \text{'MI'}) \}$ \\
\addlinespace[3pt]
$\square$\ \textbf{(B)} $\{ p.pname \mid p \in teams \land \exists t (t \in players \land p.tid = t.tid \land t.tname = \text{'MI'}) \}$ \\
\addlinespace[3pt]
$\square$\ \textbf{(C)} $\{ p.pname \mid p \in players \land \exists t (t \in teams \land t.tname = \text{'MI'}) \}$ \\
\addlinespace[3pt]
$\square$\ \textbf{(D)} $\{ p.pname \mid p \in teams \land \exists t (t \in players \land t.tname = \text{'MI'}) \}$
\end{tabular}

\tcblower

\textbf{Official Key:} \textbf{\large A} \hfill \textit{\small [Medium]}
\vspace{2pt}\hrule height 0.3pt\vspace{3pt}
\textbf{Topic:} Databases \\
\textbf{Subtopic:} \textit{Relational model: relational algebra, tuple calculus, SQL}
\vspace{2pt}\hrule height 0.3pt\vspace{3pt}
\textbf{Guide / Hint:} \\
{\small The query must select `p.pname` from the `players` relation ($p \in players$) and join with a tuple $t \in teams$ on the foreign key condition $p.tid = t.tid$, while filtering for $t.tname = \text{'MI'}$. Option A accurately expresses this quantified condition in Tuple Relational Calculus.}

\end{tcolorbox}
\vspace{6pt}\hrule height 0.35pt\vspace{6pt}
\noindent
\begin{tcolorbox}[
    blanker,
    sidebyside,
    sidebyside align=top,
    lefthand width=0.55\textwidth,
    top=4pt, bottom=4pt, left=0pt, right=0pt,
    middle=16pt,
    underlay={
        \draw[dashed, black!80, line width=0.8pt, dash pattern=on 4pt off 3pt]
            ([xshift=0.55\textwidth+8pt]frame.north west) --
            ([xshift=0.55\textwidth+8pt]frame.south west);
    }
]
\textbf{Q.40} \quad {\footnotesize\textsc{[MCQ $\cdot$ 2 Mark(s) $\cdot$ Neg: 0.66]}} \hfill {\tiny\texttt{2025\_CS1\_CS\_Q40}}
\vspace{2pt}\hrule height 0.3pt\vspace{4pt}

A packet with the destination IP address 145.36.109.70 arrives at a router whose routing table is shown below:
\par\medskip
\begin{tabular}{|l|l|c|}
\hline
\textbf{Subnet Address} & \textbf{Subnet Mask (CIDR)} & \textbf{Interface} \\
\hline
145.36.0.0 & /16 & E1 \\
145.36.128.0 & /17 & E2 \\
145.36.64.0 & /18 & E3 \\
145.36.255.0 & /24 & E4 \\
Default & -- & E5 \\
\hline
\end{tabular}
\par\medskip
Which interface will the packet be forwarded to?

\par\vspace{6pt}\noindent
\begin{tabular}{@{}p{0.48\linewidth}@{\hspace{0.04\linewidth}}p{0.48\linewidth}@{}}
$\square$\ \textbf{(A)} E3 & $\square$\ \textbf{(B)} E1 \\
\addlinespace[3pt]
$\square$\ \textbf{(C)} E2 & $\square$\ \textbf{(D)} E5
\end{tabular}

\tcblower

\textbf{Official Key:} \textbf{\large A} \hfill \textit{\small [Medium]}
\vspace{2pt}\hrule height 0.3pt\vspace{3pt}
\textbf{Topic:} Computer Networks \\
\textbf{Subtopic:} \textit{Routing protocols: shortest path, flooding, distance vector and link state routing}
\vspace{2pt}\hrule height 0.3pt\vspace{3pt}
\textbf{Guide / Hint:} \\
{\small The destination address has third octet $109 = 01101101_2$. For /16, $145.36.0.0$ matches. For /17 ($128 = 10000000_2$), the leading bit is 0, so /17 does not match. For /18 ($64 = 01000000_2$), the first two bits are $01_2$, which matches $109$ ($01101101_2$). For /24, 255 does not match. By Longest Prefix Match, /18 (interface E3) is chosen.}

\end{tcolorbox}
\vspace{6pt}\hrule height 0.35pt\vspace{6pt}
\noindent
\begin{tcolorbox}[
    blanker,
    sidebyside,
    sidebyside align=top,
    lefthand width=0.55\textwidth,
    top=4pt, bottom=4pt, left=0pt, right=0pt,
    middle=16pt,
    underlay={
        \draw[dashed, black!80, line width=0.8pt, dash pattern=on 4pt off 3pt]
            ([xshift=0.55\textwidth+8pt]frame.north west) --
            ([xshift=0.55\textwidth+8pt]frame.south west);
    }
]
\textbf{Q.41} \quad {\footnotesize\textsc{[MCQ $\cdot$ 2 Mark(s) $\cdot$ Neg: 0.66]}} \hfill {\tiny\texttt{2025\_CS1\_CS\_Q41}}
\vspace{2pt}\hrule height 0.3pt\vspace{4pt}

Let $A$ be a $2 \times 2$ matrix given by:
\[ A = \begin{bmatrix} 1 & 1 \\ 1 & -1 \end{bmatrix} \]
What are the eigenvalues of the matrix $A^{13}$?

\par\vspace{6pt}\noindent
\begin{tabular}{@{}p{0.48\linewidth}@{\hspace{0.04\linewidth}}p{0.48\linewidth}@{}}
$\square$\ \textbf{(A)} $1, -1$ & $\square$\ \textbf{(B)} $2\sqrt{2}, -2\sqrt{2}$ \\
\addlinespace[3pt]
$\square$\ \textbf{(C)} $4\sqrt{2}, -4\sqrt{2}$ & $\square$\ \textbf{(D)} $64\sqrt{2}, -64\sqrt{2}$
\end{tabular}

\tcblower

\textbf{Official Key:} \textbf{\large D} \hfill \textit{\small [Medium]}
\vspace{2pt}\hrule height 0.3pt\vspace{3pt}
\textbf{Topic:} Linear Algebra \\
\textbf{Subtopic:} \textit{Eigenvalues and eigenvectors}
\vspace{2pt}\hrule height 0.3pt\vspace{3pt}
\textbf{Guide / Hint:} \\
{\small The characteristic equation of $A$ is $\det(A - \lambda I) = (1 - \lambda)(-1 - \lambda) - 1 = \lambda^2 - 2 = 0$, so eigenvalues of $A$ are $\lambda_1 = \sqrt{2}$ and $\lambda_2 = -\sqrt{2}$. The eigenvalues of $A^{13}$ are $(\sqrt{2})^{13} = 2^6 \sqrt{2} = 64\sqrt{2}$ and $(-\sqrt{2})^{13} = -64\sqrt{2}$.}

\end{tcolorbox}
\vspace{6pt}\hrule height 0.35pt\vspace{6pt}
\noindent
\begin{tcolorbox}[
    blanker,
    sidebyside,
    sidebyside align=top,
    lefthand width=0.55\textwidth,
    top=4pt, bottom=4pt, left=0pt, right=0pt,
    middle=16pt,
    underlay={
        \draw[dashed, black!80, line width=0.8pt, dash pattern=on 4pt off 3pt]
            ([xshift=0.55\textwidth+8pt]frame.north west) --
            ([xshift=0.55\textwidth+8pt]frame.south west);
    }
]
\textbf{Q.42} \quad {\footnotesize\textsc{[MCQ $\cdot$ 2 Mark(s) $\cdot$ Neg: 0.66]}} \hfill {\tiny\texttt{2025\_CS1\_CS\_Q42}}
\vspace{2pt}\hrule height 0.3pt\vspace{4pt}

Consider the following four variable Boolean function in sum-of-product form:
\[ F(b\_3, b\_2, b\_1, b\_0) = \sum(0, 2, 4, 8, 10, 11, 12) \]
where the value of the function is computed by considering $b_3b_2b_1b_0$ as a 4-bit binary number, where $b_3$ denotes the most significant bit and $b_0$ denotes the least significant bit. Note that there are no don't care terms. Which ONE of the following options is the CORRECT minimized expression?

\par\vspace{6pt}\noindent
\begin{tabular}{@{}p{0.98\linewidth}@{}}
$\square$\ \textbf{(A)} $\overline{b}_1\overline{b}_0 + \overline{b}_2\overline{b}_0 + b_1\overline{b}_2b_3$ \\
\addlinespace[3pt]
$\square$\ \textbf{(B)} $\overline{b}_1\overline{b}_0 + \overline{b}_2\overline{b}_0$ \\
\addlinespace[3pt]
$\square$\ \textbf{(C)} $\overline{b}_2\overline{b}_0 + b_1b_2b_3$ \\
\addlinespace[3pt]
$\square$\ \textbf{(D)} $\overline{b}_0\overline{b}_2 + \overline{b}_3$
\end{tabular}

\tcblower

\textbf{Official Key:} \textbf{\large A} \hfill \textit{\small [Medium]}
\vspace{2pt}\hrule height 0.3pt\vspace{3pt}
\textbf{Topic:} Digital Logic \\
\textbf{Subtopic:} \textit{Minimization}
\vspace{2pt}\hrule height 0.3pt\vspace{3pt}
\textbf{Guide / Hint:} \\
{\small Grouping minterms on a 4-variable K-map: minterms $\{0, 2, 8, 10\}$ form a 4-cell group giving $\overline{b}_2\overline{b}_0$; minterms $\{0, 4, 8, 12\}$ form a 4-cell group giving $\overline{b}_1\overline{b}_0$; remaining minterm 11 combines with 10 to give $b_3\overline{b}_2b_1$. Combining these essential prime implicants yields $\overline{b}_1\overline{b}_0 + \overline{b}_2\overline{b}_0 + b_1\overline{b}_2b_3$.}

\end{tcolorbox}
\vspace{6pt}\hrule height 0.35pt\vspace{6pt}
\noindent
\begin{tcolorbox}[
    blanker,
    sidebyside,
    sidebyside align=top,
    lefthand width=0.55\textwidth,
    top=4pt, bottom=4pt, left=0pt, right=0pt,
    middle=16pt,
    underlay={
        \draw[dashed, black!80, line width=0.8pt, dash pattern=on 4pt off 3pt]
            ([xshift=0.55\textwidth+8pt]frame.north west) --
            ([xshift=0.55\textwidth+8pt]frame.south west);
    }
]
\textbf{Q.43} \quad {\footnotesize\textsc{[MCQ $\cdot$ 2 Mark(s) $\cdot$ Neg: 0.66]}} \hfill {\tiny\texttt{2025\_CS1\_CS\_Q43}}
\vspace{2pt}\hrule height 0.3pt\vspace{4pt}

Let $G(V, E)$ be an undirected and unweighted graph with 100 vertices. Let $d(u, v)$ denote the number of edges in a shortest path between vertices $u$ and $v$ in $V$. Let the maximum value of $d(u, v), u, v \in V$ such that $u \neq v$, be 30. Let $T$ be any breadth-first-search tree of $G$. Which ONE of the given options is CORRECT for every such graph $G$?

\par\vspace{6pt}\noindent
\begin{tabular}{@{}p{0.98\linewidth}@{}}
$\square$\ \textbf{(A)} The height of $T$ is exactly 15. \\
\addlinespace[3pt]
$\square$\ \textbf{(B)} The height of $T$ is exactly 30. \\
\addlinespace[3pt]
$\square$\ \textbf{(C)} The height of $T$ is at least 15. \\
\addlinespace[3pt]
$\square$\ \textbf{(D)} The height of $T$ is at least 30.
\end{tabular}

\tcblower

\textbf{Official Key:} \textbf{\large C} \hfill \textit{\small [Medium]}
\vspace{2pt}\hrule height 0.3pt\vspace{3pt}
\textbf{Topic:} Algorithms \\
\textbf{Subtopic:} \textit{Graph traversals, minimum spanning trees, shortest paths}
\vspace{2pt}\hrule height 0.3pt\vspace{3pt}
\textbf{Guide / Hint:} \\
{\small The diameter of $G$ is $\text{diam}(G) = \max_{u,v} d(u, v) = 30$. For any BFS root $r$, the triangle inequality ensures that for the pair $(u, v)$ realizing the diameter, $30 = d(u, v) \le d(r, u) + d(r, v) \le 2 \times \text{height}(T)$. Hence, $\text{height}(T) \ge \lceil 30/2 \rceil = 15$.}

\end{tcolorbox}
\vspace{6pt}\hrule height 0.35pt\vspace{6pt}
\noindent
\begin{tcolorbox}[
    blanker,
    sidebyside,
    sidebyside align=top,
    lefthand width=0.55\textwidth,
    top=4pt, bottom=4pt, left=0pt, right=0pt,
    middle=16pt,
    underlay={
        \draw[dashed, black!80, line width=0.8pt, dash pattern=on 4pt off 3pt]
            ([xshift=0.55\textwidth+8pt]frame.north west) --
            ([xshift=0.55\textwidth+8pt]frame.south west);
    }
]
\textbf{Q.44} \quad {\footnotesize\textsc{[MCQ $\cdot$ 2 Mark(s) $\cdot$ Neg: 0.66]}} \hfill {\tiny\texttt{2025\_CS1\_CS\_Q44}}
\vspace{2pt}\hrule height 0.3pt\vspace{4pt}

Consider the following two languages over the alphabet $\{a, b\}$:
\[ L\_1 = \{ \alpha\beta\alpha \mid \alpha \in \{a, b\}^+ \text{ AND } \beta \in \{a, b\}^+ \} \]
\[ L\_2 = \{ \alpha\beta\alpha \mid \alpha \in \{a\}^+ \text{ AND } \beta \in \{a, b\}^+ \} \]
Which ONE of the following statements is CORRECT?

\par\vspace{6pt}\noindent
\begin{tabular}{@{}p{0.98\linewidth}@{}}
$\square$\ \textbf{(A)} Both $L_1$ and $L_2$ are regular languages. \\
\addlinespace[3pt]
$\square$\ \textbf{(B)} $L_1$ is a regular language but $L_2$ is not a regular language. \\
\addlinespace[3pt]
$\square$\ \textbf{(C)} $L_1$ is not a regular language but $L_2$ is a regular language. \\
\addlinespace[3pt]
$\square$\ \textbf{(D)} Neither $L_1$ nor $L_2$ is a regular language.
\end{tabular}

\tcblower

\textbf{Official Key:} \textbf{\large C} \hfill \textit{\small [Medium]}
\vspace{2pt}\hrule height 0.3pt\vspace{3pt}
\textbf{Topic:} Theory of Computation \\
\textbf{Subtopic:} \textit{Regular and context-free languages, pumping lemma}
\vspace{2pt}\hrule height 0.3pt\vspace{3pt}
\textbf{Guide / Hint:} \\
{\small For $L_2$, $\alpha = a^k$ ($k \ge 1$). A string belongs to $L_2$ if and only if it starts with $a$, ends with $a$, and has length at least 3 ($a(a+b)^+a$), which is clearly regular. For $L_1$, the prefix $\alpha$ and suffix $\alpha$ must match over $\{a, b\}^+$, which requires unbounded memory to verify matching arbitrary binary strings, making $L_1$ non-regular.}

\end{tcolorbox}
\vspace{6pt}\hrule height 0.35pt\vspace{6pt}
\noindent
\begin{tcolorbox}[
    blanker,
    sidebyside,
    sidebyside align=top,
    lefthand width=0.55\textwidth,
    top=4pt, bottom=4pt, left=0pt, right=0pt,
    middle=16pt,
    underlay={
        \draw[dashed, black!80, line width=0.8pt, dash pattern=on 4pt off 3pt]
            ([xshift=0.55\textwidth+8pt]frame.north west) --
            ([xshift=0.55\textwidth+8pt]frame.south west);
    }
]
\textbf{Q.45} \quad {\footnotesize\textsc{[MCQ $\cdot$ 2 Mark(s) $\cdot$ Neg: 0.66]}} \hfill {\tiny\texttt{2025\_CS1\_CS\_Q45}}
\vspace{2pt}\hrule height 0.3pt\vspace{4pt}

Consider the following two languages over the alphabet $\{a, b, c\}$, where $m$ and $n$ are natural numbers.
\[ L\_1 = \{a^mb^mc^{m+n} \mid m, n \ge 1\} \]
\[ L\_2 = \{a^mb^nc^{m+n} \mid m, n \ge 1\} \]
Which ONE of the following statements is CORRECT?

\par\vspace{6pt}\noindent
\begin{tabular}{@{}p{0.98\linewidth}@{}}
$\square$\ \textbf{(A)} Both $L_1$ and $L_2$ are context-free languages. \\
\addlinespace[3pt]
$\square$\ \textbf{(B)} $L_1$ is a context-free language but $L_2$ is not a context-free language. \\
\addlinespace[3pt]
$\square$\ \textbf{(C)} $L_1$ is not a context-free language but $L_2$ is a context-free language. \\
\addlinespace[3pt]
$\square$\ \textbf{(D)} Neither $L_1$ nor $L_2$ are context-free languages.
\end{tabular}

\tcblower

\textbf{Official Key:} \textbf{\large C} \hfill \textit{\small [Medium]}
\vspace{2pt}\hrule height 0.3pt\vspace{3pt}
\textbf{Topic:} Theory of Computation \\
\textbf{Subtopic:} \textit{Context-free grammars and push-down automata}
\vspace{2pt}\hrule height 0.3pt\vspace{3pt}
\textbf{Guide / Hint:} \\
{\small In $L_2$, rewrite $c^{m+n} = c^n c^m$ or generate using PDA by pushing $a$'s and $b$'s and popping on $c$'s ($c$'s match total count of $a$'s and $b$'s), which requires a single stack, so $L_2$ is context-free. In $L_1$, $a^m b^m c^m c^n$ simultaneously requires equality between powers of $a$ and $b$, plus matching with $c$, requiring two independent stack comparisons, making $L_1$ non-CFL.}

\end{tcolorbox}
\vspace{6pt}\hrule height 0.35pt\vspace{6pt}
\noindent
\begin{tcolorbox}[
    blanker,
    sidebyside,
    sidebyside align=top,
    lefthand width=0.55\textwidth,
    top=4pt, bottom=4pt, left=0pt, right=0pt,
    middle=16pt,
    underlay={
        \draw[dashed, black!80, line width=0.8pt, dash pattern=on 4pt off 3pt]
            ([xshift=0.55\textwidth+8pt]frame.north west) --
            ([xshift=0.55\textwidth+8pt]frame.south west);
    }
]
\textbf{Q.46} \quad {\footnotesize\textsc{[MSQ $\cdot$ 2 Mark(s) $\cdot$ No Neg]}} \hfill {\tiny\texttt{2025\_CS1\_CS\_Q46}}
\vspace{2pt}\hrule height 0.3pt\vspace{4pt}

Which of the following statement(s) is/are TRUE while computing First and Follow during top down parsing by a compiler?

\par\vspace{6pt}\noindent
\begin{tabular}{@{}p{0.98\linewidth}@{}}
$\square$\ \textbf{(A)} For a production $A \to \epsilon$, $\epsilon$ will be added to $First(A)$. \\
\addlinespace[3pt]
$\square$\ \textbf{(B)} If there is any input right end marker, it will be added to $First(S)$, where $S$ is the start symbol. \\
\addlinespace[3pt]
$\square$\ \textbf{(C)} For a production $A \to \epsilon$, $\epsilon$ will be added to $Follow(A)$. \\
\addlinespace[3pt]
$\square$\ \textbf{(D)} If there is any input right end marker, it will be added to $Follow(S)$, where $S$ is the start symbol.
\end{tabular}

\tcblower

\textbf{Official Key:} \textbf{\large A;D} \hfill \textit{\small [Easy]}
\vspace{2pt}\hrule height 0.3pt\vspace{3pt}
\textbf{Topic:} Compiler Design \\
\textbf{Subtopic:} \textit{Lexical analysis, parsing, syntax-directed translation}
\vspace{2pt}\hrule height 0.3pt\vspace{3pt}
\textbf{Guide / Hint:} \\
{\small By definition of FIRST and FOLLOW: if $A \to \epsilon$, then $\epsilon \in \text{FIRST}(A)$ (A is true). $\epsilon$ can never belong to any FOLLOW set (C is false). The right end marker (\$) is placed into $\text{FOLLOW}(S)$ for the start symbol $S$ to mark the end of input (D is true, B is false).}

\end{tcolorbox}
\vspace{6pt}\hrule height 0.35pt\vspace{6pt}
\noindent
\begin{tcolorbox}[
    blanker,
    sidebyside,
    sidebyside align=top,
    lefthand width=0.55\textwidth,
    top=4pt, bottom=4pt, left=0pt, right=0pt,
    middle=16pt,
    underlay={
        \draw[dashed, black!80, line width=0.8pt, dash pattern=on 4pt off 3pt]
            ([xshift=0.55\textwidth+8pt]frame.north west) --
            ([xshift=0.55\textwidth+8pt]frame.south west);
    }
]
\textbf{Q.47} \quad {\footnotesize\textsc{[MSQ $\cdot$ 2 Mark(s) $\cdot$ No Neg]}} \hfill {\tiny\texttt{2025\_CS1\_CS\_Q47}}
\vspace{2pt}\hrule height 0.3pt\vspace{4pt}

Consider a relational schema $team(name, city, owner)$, with functional dependencies $\{name \to city, name \to owner\}$. The relation $team$ is decomposed into two relations, $t_1(name, city)$ and $t_2(name, owner)$. Which of the following statement(s) is/are TRUE?

\par\vspace{6pt}\noindent
\begin{tabular}{@{}p{0.98\linewidth}@{}}
$\square$\ \textbf{(A)} The relation $team$ is NOT in BCNF. \\
\addlinespace[3pt]
$\square$\ \textbf{(B)} The relations $t_1$ and $t_2$ are in BCNF. \\
\addlinespace[3pt]
$\square$\ \textbf{(C)} The decomposition constitutes a lossless join. \\
\addlinespace[3pt]
$\square$\ \textbf{(D)} The relation $team$ is NOT in 3NF.
\end{tabular}

\tcblower

\textbf{Official Key:} \textbf{\large B;C} \hfill \textit{\small [Medium]}
\vspace{2pt}\hrule height 0.3pt\vspace{3pt}
\textbf{Topic:} Databases \\
\textbf{Subtopic:} \textit{Integrity constraints, normal forms}
\vspace{2pt}\hrule height 0.3pt\vspace{3pt}
\textbf{Guide / Hint:} \\
{\small In $team$, $name$ is a candidate key since $name^+ = \{name, city, owner\}$. Since every FD has a superkey on the LHS, $team$ was already in BCNF and 3NF (A and D are false). In $t_1$ and $t_2$, $name$ is the primary key, so both are in BCNF (B is true). The common attribute $R_1 \cap R_2 = \{name\}$ is a key of both $t_1$ and $t_2$, ensuring a lossless join (C is true).}

\end{tcolorbox}
\vspace{6pt}\hrule height 0.35pt\vspace{6pt}
\noindent
\begin{tcolorbox}[
    blanker,
    sidebyside,
    sidebyside align=top,
    lefthand width=0.55\textwidth,
    top=4pt, bottom=4pt, left=0pt, right=0pt,
    middle=16pt,
    underlay={
        \draw[dashed, black!80, line width=0.8pt, dash pattern=on 4pt off 3pt]
            ([xshift=0.55\textwidth+8pt]frame.north west) --
            ([xshift=0.55\textwidth+8pt]frame.south west);
    }
]
\textbf{Q.48} \quad {\footnotesize\textsc{[MSQ $\cdot$ 2 Mark(s) $\cdot$ No Neg]}} \hfill {\tiny\texttt{2025\_CS1\_CS\_Q48}}
\vspace{2pt}\hrule height 0.3pt\vspace{4pt}

Which of the following predicate logic formulae/formula is/are CORRECT representation(s) of the statement: ``Everyone has exactly one mother''?
The meanings of the predicates used are:
\begin{itemize}
\item $mother(y, x)$: $y$ is the mother of $x$
\item $noteq(x, y)$: $x$ and $y$ are not equal
\end{itemize}

\par\vspace{6pt}\noindent
\begin{tabular}{@{}p{0.98\linewidth}@{}}
$\square$\ \textbf{(A)} $\forall x \exists y \exists z (mother(y, x) \land \neg mother(z, x))$ \\
\addlinespace[3pt]
$\square$\ \textbf{(B)} $\forall x \exists y [mother(y, x) \land \forall z (noteq(z, y) \to \neg mother(z, x))]$ \\
\addlinespace[3pt]
$\square$\ \textbf{(C)} $\forall x \forall y [mother(y, x) \to \exists z (mother(z, x) \land \neg noteq(z, y))]$ \\
\addlinespace[3pt]
$\square$\ \textbf{(D)} $\forall x \exists y [mother(y, x) \land \neg \exists z (noteq(z, y) \land mother(z, x))]$
\end{tabular}

\tcblower

\textbf{Official Key:} \textbf{\large B;D} \hfill \textit{\small [Medium]}
\vspace{2pt}\hrule height 0.3pt\vspace{3pt}
\textbf{Topic:} Discrete Mathematics \\
\textbf{Subtopic:} \textit{Propositional and first order logic}
\vspace{2pt}\hrule height 0.3pt\vspace{3pt}
\textbf{Guide / Hint:} \\
{\small ``Exactly one mother'' requires existence ($\exists y \, mother(y, x)$) and uniqueness (any $z \neq y$ cannot be a mother of $x$). Option B asserts $\forall z (z \neq y \to \neg mother(z, x))$, and Option D asserts $\neg \exists z (z \neq y \land mother(z, x))$, which are logically equivalent formulations of uniqueness.}

\end{tcolorbox}
\vspace{6pt}\hrule height 0.35pt\vspace{6pt}
\noindent
\begin{tcolorbox}[
    blanker,
    sidebyside,
    sidebyside align=top,
    lefthand width=0.55\textwidth,
    top=4pt, bottom=4pt, left=0pt, right=0pt,
    middle=16pt,
    underlay={
        \draw[dashed, black!80, line width=0.8pt, dash pattern=on 4pt off 3pt]
            ([xshift=0.55\textwidth+8pt]frame.north west) --
            ([xshift=0.55\textwidth+8pt]frame.south west);
    }
]
\textbf{Q.49} \quad {\footnotesize\textsc{[MSQ $\cdot$ 2 Mark(s) $\cdot$ No Neg]}} \hfill {\tiny\texttt{2025\_CS1\_CS\_Q49}}
\vspace{2pt}\hrule height 0.3pt\vspace{4pt}

$A = \{0, 1, 2, 3, \dots \}$ is the set of non-negative integers. Let $\mathcal{F}$ be the set of functions from $A$ to itself. For any two functions $f_1, f_2 \in \mathcal{F}$, we define
\[ (f\_1 \odot f\_2)(n) = f\_1(n) + f\_2(n) \]
for every number $n$ in $A$. Which of the following is/are CORRECT about the mathematical structure $(\mathcal{F}, \odot)$?

\par\vspace{6pt}\noindent
\begin{tabular}{@{}p{0.98\linewidth}@{}}
$\square$\ \textbf{(A)} $(\mathcal{F}, \odot)$ is an Abelian group. \\
\addlinespace[3pt]
$\square$\ \textbf{(B)} $(\mathcal{F}, \odot)$ is an Abelian monoid. \\
\addlinespace[3pt]
$\square$\ \textbf{(C)} $(\mathcal{F}, \odot)$ is a non-Abelian group. \\
\addlinespace[3pt]
$\square$\ \textbf{(D)} $(\mathcal{F}, \odot)$ is a non-Abelian monoid.
\end{tabular}

\tcblower

\textbf{Official Key:} \textbf{\large B} \hfill \textit{\small [Medium]}
\vspace{2pt}\hrule height 0.3pt\vspace{3pt}
\textbf{Topic:} Discrete Mathematics \\
\textbf{Subtopic:} \textit{Monoids, Groups}
\vspace{2pt}\hrule height 0.3pt\vspace{3pt}
\textbf{Guide / Hint:} \\
{\small Addition on non-negative integers is closed, associative, and commutative, so $\odot$ is associative and commutative. The identity function is $e(n) = 0$ for all $n \in A$, making $(\mathcal{F}, \odot)$ a commutative (Abelian) monoid. Inverses do not exist because negative function values are not in $A$, so it is not a group.}

\end{tcolorbox}
\vspace{6pt}\hrule height 0.35pt\vspace{6pt}
\noindent
\begin{tcolorbox}[
    blanker,
    sidebyside,
    sidebyside align=top,
    lefthand width=0.55\textwidth,
    top=4pt, bottom=4pt, left=0pt, right=0pt,
    middle=16pt,
    underlay={
        \draw[dashed, black!80, line width=0.8pt, dash pattern=on 4pt off 3pt]
            ([xshift=0.55\textwidth+8pt]frame.north west) --
            ([xshift=0.55\textwidth+8pt]frame.south west);
    }
]
\textbf{Q.50} \quad {\footnotesize\textsc{[MSQ $\cdot$ 2 Mark(s) $\cdot$ No Neg]}} \hfill {\tiny\texttt{2025\_CS1\_CS\_Q50}}
\vspace{2pt}\hrule height 0.3pt\vspace{4pt}

Consider the deterministic finite automaton (DFA) defined over the alphabet $\Sigma = \{a, b\}$. Identify which of the following language(s) is/are accepted by the given DFA.

\par\vspace{6pt}\noindent
\begin{tabular}{@{}p{0.98\linewidth}@{}}
$\square$\ \textbf{(A)} The set of all strings containing an even number of $b$'s. \\
\addlinespace[3pt]
$\square$\ \textbf{(B)} The set of all strings containing the pattern $bab$. \\
\addlinespace[3pt]
$\square$\ \textbf{(C)} The set of all strings ending with the pattern $bab$. \\
\addlinespace[3pt]
$\square$\ \textbf{(D)} The set of all strings not containing the pattern $aba$.
\end{tabular}

\tcblower

\textbf{Official Key:} \textbf{\large C} \hfill \textit{\small [Medium]}
\vspace{2pt}\hrule height 0.3pt\vspace{3pt}
\textbf{Topic:} Theory of Computation \\
\textbf{Subtopic:} \textit{Regular expressions and finite automata}
\vspace{2pt}\hrule height 0.3pt\vspace{3pt}
\textbf{Guide / Hint:} \\
{\small In the given DFA, the sequence of transitions leading to the unique accepting state $q_3$ corresponds precisely to recognizing the suffix $bab$. Any character following $q_3$ transitions back according to pattern matching prefixes of $bab$, accepting exactly strings ending with the pattern $bab$.}

\end{tcolorbox}
\vspace{6pt}\hrule height 0.35pt\vspace{6pt}
\noindent
\begin{tcolorbox}[
    blanker,
    sidebyside,
    sidebyside align=top,
    lefthand width=0.55\textwidth,
    top=4pt, bottom=4pt, left=0pt, right=0pt,
    middle=16pt,
    underlay={
        \draw[dashed, black!80, line width=0.8pt, dash pattern=on 4pt off 3pt]
            ([xshift=0.55\textwidth+8pt]frame.north west) --
            ([xshift=0.55\textwidth+8pt]frame.south west);
    }
]
\textbf{Q.51} \quad {\footnotesize\textsc{[NAT $\cdot$ 2 Mark(s) $\cdot$ No Neg]}} \hfill {\tiny\texttt{2025\_CS1\_CS\_Q51}}
\vspace{2pt}\hrule height 0.3pt\vspace{4pt}

A disk of size 512M bytes is divided into blocks of 64K bytes. A file is stored in the disk using linked allocation. In linked allocation, each data block reserves 4 bytes to store the pointer to the next data block. The link part of the last data block contains a NULL pointer (also of 4 bytes). Suppose a file of 1M bytes needs to be stored in the disk. Assume $1\text{K} = 2^{10}$ and $1\text{M} = 2^{20}$. The amount of space in bytes that will be wasted due to internal fragmentation is \_\_\_\_\_\_\. (Answer in integer)

\par\vspace{6pt}\noindent\textbf{Your Answer:} \fbox{\rule{0pt}{1.3ex}\hspace{3cm}}

\tcblower

\textbf{Official Key:} \textbf{\large 65468 to 65468} \hfill \textit{\small [Medium]}
\vspace{2pt}\hrule height 0.3pt\vspace{3pt}
\textbf{Topic:} Operating System \\
\textbf{Subtopic:} \textit{File systems}
\vspace{2pt}\hrule height 0.3pt\vspace{3pt}
\textbf{Guide / Hint:} \\
{\small Block size is $64\text{ KB} = 65536$ bytes, with 4 bytes for pointer, leaving $65536 - 4 = 65532$ bytes of file data per block. File size is $1\text{ MB} = 1048576$ bytes. Number of blocks needed is $\lceil 1048576 / 65532 \rceil = \lceil 16.001037 \rceil = 17$ blocks. The first 16 blocks store $16 \times 65532 = 1048512$ bytes, leaving $1048576 - 1048512 = 64$ bytes stored in the 17th block. Internal fragmentation in the last block is $65532 - 64 = 65468$ bytes.}

\end{tcolorbox}
\vspace{6pt}\hrule height 0.35pt\vspace{6pt}
\noindent
\begin{tcolorbox}[
    blanker,
    sidebyside,
    sidebyside align=top,
    lefthand width=0.55\textwidth,
    top=4pt, bottom=4pt, left=0pt, right=0pt,
    middle=16pt,
    underlay={
        \draw[dashed, black!80, line width=0.8pt, dash pattern=on 4pt off 3pt]
            ([xshift=0.55\textwidth+8pt]frame.north west) --
            ([xshift=0.55\textwidth+8pt]frame.south west);
    }
]
\textbf{Q.52} \quad {\footnotesize\textsc{[NAT $\cdot$ 2 Mark(s) $\cdot$ No Neg]}} \hfill {\tiny\texttt{2025\_CS1\_CS\_Q52}}
\vspace{2pt}\hrule height 0.3pt\vspace{4pt}

Refer to the given 3-address code sequence. This code sequence is split into basic blocks. The number of basic blocks is \_\_\_\_\_\_\. (Answer in integer)
\par\medskip
\begin{verbatim}
1001: i = 1
1002: j = 1
1003: t1 = 10*i
1004: t2 = t1+j
1005: t3 = 8*t2
1006: t4 = t3-88
1007: a[t4] = 0.0
1008: j = j+1
1009: if j <= 10 goto 1003
1010: i = i+1
1011: if i <= 10 goto 1002
1012: i = 1
1013: t5 = i-1
1014: t6 = 88*t5
1015: a[t6] = 1.0
1016: i = i+1
1017: if i <= 11 goto 1013
\end{verbatim}

\par\vspace{6pt}\noindent\textbf{Your Answer:} \fbox{\rule{0pt}{1.3ex}\hspace{3cm}}

\tcblower

\textbf{Official Key:} \textbf{\large 6 to 6} \hfill \textit{\small [Medium]}
\vspace{2pt}\hrule height 0.3pt\vspace{3pt}
\textbf{Topic:} Compiler Design \\
\textbf{Subtopic:} \textit{Intermediate code generation}
\vspace{2pt}\hrule height 0.3pt\vspace{3pt}
\textbf{Guide / Hint:} \\
{\small Leaders are: 1001 (first instruction), targets of jumps: 1003, 1002, 1013, and instructions immediately following jumps: 1010, 1012. Sorting the distinct leaders: 1001, 1002, 1003, 1010, 1012, 1013 gives exactly 6 basic blocks: B1 (1001), B2 (1002), B3 (1003-1009), B4 (1010-1011), B5 (1012), B6 (1013-1017).}

\end{tcolorbox}
\vspace{6pt}\hrule height 0.35pt\vspace{6pt}
\noindent
\begin{tcolorbox}[
    blanker,
    sidebyside,
    sidebyside align=top,
    lefthand width=0.55\textwidth,
    top=4pt, bottom=4pt, left=0pt, right=0pt,
    middle=16pt,
    underlay={
        \draw[dashed, black!80, line width=0.8pt, dash pattern=on 4pt off 3pt]
            ([xshift=0.55\textwidth+8pt]frame.north west) --
            ([xshift=0.55\textwidth+8pt]frame.south west);
    }
]
\textbf{Q.53} \quad {\footnotesize\textsc{[NAT $\cdot$ 2 Mark(s) $\cdot$ No Neg]}} \hfill {\tiny\texttt{2025\_CS1\_CS\_Q53}}
\vspace{2pt}\hrule height 0.3pt\vspace{4pt}

A computer has a memory hierarchy consisting of two-level cache (L1 and L2) and a main memory. If the processor needs to access data from memory, it first looks into L1 cache ($h_1 = 0.95$, access time $t_1 = 10\text{ ns}$). If not found in L1, it accesses L2 cache ($h_2 = 0.85$, access time including L1 miss penalty $t_2 = 20\text{ ns}$). If it misses in L2, it accesses main memory (access time including L1 and L2 miss penalties $t_{MM} = 200\text{ ns}$). The average memory access time in nanoseconds (ns) is \_\_\_\_\_\_\. (rounded off to two decimal places)

\par\vspace{6pt}\noindent\textbf{Your Answer:} \fbox{\rule{0pt}{1.3ex}\hspace{3cm}}

\tcblower

\textbf{Official Key:} \textbf{\large 11.83 to 11.87} \hfill \textit{\small [Medium]}
\vspace{2pt}\hrule height 0.3pt\vspace{3pt}
\textbf{Topic:} Computer Organization and Architecture \\
\textbf{Subtopic:} \textit{Memory hierarchy: cache, main memory and secondary storage}
\vspace{2pt}\hrule height 0.3pt\vspace{3pt}
\textbf{Guide / Hint:} \\
{\small Using simultaneous access with cumulative times: $\text{AMAT} = h_1 t_1 + (1 - h_1) h_2 t_2 + (1 - h_1)(1 - h_2) t_{MM} = 0.95(10) + 0.05(0.85)(20) + 0.05(0.15)(200) = 9.5 + 0.85 + 1.5 = 11.85\text{ ns}$.}

\end{tcolorbox}
\vspace{6pt}\hrule height 0.35pt\vspace{6pt}
\noindent
\begin{tcolorbox}[
    blanker,
    sidebyside,
    sidebyside align=top,
    lefthand width=0.55\textwidth,
    top=4pt, bottom=4pt, left=0pt, right=0pt,
    middle=16pt,
    underlay={
        \draw[dashed, black!80, line width=0.8pt, dash pattern=on 4pt off 3pt]
            ([xshift=0.55\textwidth+8pt]frame.north west) --
            ([xshift=0.55\textwidth+8pt]frame.south west);
    }
]
\textbf{Q.54} \quad {\footnotesize\textsc{[NAT $\cdot$ 2 Mark(s) $\cdot$ No Neg]}} \hfill {\tiny\texttt{2025\_CS1\_CS\_Q54}}
\vspace{2pt}\hrule height 0.3pt\vspace{4pt}

In optimal page replacement algorithm, information about all future page references is available to the operating system (OS). A modification of the optimal page replacement algorithm is as follows:
The OS correctly predicts only up to next 4 page references (including the current page) at the time of allocating a frame to a page.
A process accesses the pages in the following order of page numbers:
1, 3, 2, 4, 2, 3, 1, 2, 4, 3, 1, 4.
If the system has three memory frames that are initially empty, the number of page faults that will occur during execution of the process is \_\_\_\_\_\_\. (Answer in integer)

\par\vspace{6pt}\noindent\textbf{Your Answer:} \fbox{\rule{0pt}{1.3ex}\hspace{3cm}}

\tcblower

\textbf{Official Key:} \textbf{\large 6 to 6} \hfill \textit{\small [Medium]}
\vspace{2pt}\hrule height 0.3pt\vspace{3pt}
\textbf{Topic:} Operating System \\
\textbf{Subtopic:} \textit{Memory management and virtual memory}
\vspace{2pt}\hrule height 0.3pt\vspace{3pt}
\textbf{Guide / Hint:} \\
{\small Simulating the 3 frames with 4-step lookahead: 1 (fault, [1]), 3 (fault, [1, 3]), 2 (fault, [1, 3, 2]). At page 4 (fault): window is (4, 2, 3, 1), farthest is 1, so replace 1 -> [4, 3, 2]. Pages 2, 3 hit. At page 1 (fault): window (1, 2, 4, 3), replace 3 -> [4, 1, 2]. Page 2, 4, 3 (fault): replace 1 -> [4, 3, 2]. Total faults = 3 compulsory + 3 replacement = 6 page faults.}

\end{tcolorbox}
\vspace{6pt}\hrule height 0.35pt\vspace{6pt}
\noindent
\begin{tcolorbox}[
    blanker,
    sidebyside,
    sidebyside align=top,
    lefthand width=0.55\textwidth,
    top=4pt, bottom=4pt, left=0pt, right=0pt,
    middle=16pt,
    underlay={
        \draw[dashed, black!80, line width=0.8pt, dash pattern=on 4pt off 3pt]
            ([xshift=0.55\textwidth+8pt]frame.north west) --
            ([xshift=0.55\textwidth+8pt]frame.south west);
    }
]
\textbf{Q.56} \quad {\footnotesize\textsc{[NAT $\cdot$ 2 Mark(s) $\cdot$ No Neg]}} \hfill {\tiny\texttt{2025\_CS1\_CS\_Q56}}
\vspace{2pt}\hrule height 0.3pt\vspace{4pt}

Suppose a 5-bit message is transmitted from a source to a destination through a noisy channel. The probability that a bit of the message gets flipped during transmission is 0.01. Flipping of each bit is independent of one another. The probability that the message is delivered error-free to the destination is \_\_\_\_\_\_\. (rounded off to three decimal places)

\par\vspace{6pt}\noindent\textbf{Your Answer:} \fbox{\rule{0pt}{1.3ex}\hspace{3cm}}

\tcblower

\textbf{Official Key:} \textbf{\large 0.949 to 0.952} \hfill \textit{\small [Easy]}
\vspace{2pt}\hrule height 0.3pt\vspace{3pt}
\textbf{Topic:} Computer Networks \\
\textbf{Subtopic:} \textit{Data link layer: framing, error detection, Medium Access Control, Ethernet bridging}
\vspace{2pt}\hrule height 0.3pt\vspace{3pt}
\textbf{Guide / Hint:} \\
{\small For each bit, the probability of correct transmission without flipping is $1 - 0.01 = 0.99$. Since all 5 bit transmissions are mutually independent, the probability that all 5 bits are transmitted error-free is $(0.99)^5 = 0.95099 \approx 0.951$.}

\end{tcolorbox}
\vspace{6pt}\hrule height 0.35pt\vspace{6pt}
\noindent
\begin{tcolorbox}[
    blanker,
    sidebyside,
    sidebyside align=top,
    lefthand width=0.55\textwidth,
    top=4pt, bottom=4pt, left=0pt, right=0pt,
    middle=16pt,
    underlay={
        \draw[dashed, black!80, line width=0.8pt, dash pattern=on 4pt off 3pt]
            ([xshift=0.55\textwidth+8pt]frame.north west) --
            ([xshift=0.55\textwidth+8pt]frame.south west);
    }
]
\textbf{Q.57} \quad {\footnotesize\textsc{[NAT $\cdot$ 2 Mark(s) $\cdot$ No Neg]}} \hfill {\tiny\texttt{2025\_CS1\_CS\_Q57}}
\vspace{2pt}\hrule height 0.3pt\vspace{4pt}

Suppose a message of size 15000 bytes is transmitted from a source to a destination using IPv4 protocol via two routers as shown: Source $\to$ Router-1 ($\text{MTU} = 5000$ bytes) $\to$ Router-2 ($\text{MTU} = 3000$ bytes) $\to$ Destination. Each router has a defined maximum transmission unit (MTU) including the 20-byte IP header. The number of fragments that will be delivered to the destination is \_\_\_\_\_\_\. (Answer in integer)

\par\vspace{6pt}\noindent\textbf{Your Answer:} \fbox{\rule{0pt}{1.3ex}\hspace{3cm}}

\tcblower

\textbf{Official Key:} \textbf{\large 7 to 7} \hfill \textit{\small [Hard]}
\vspace{2pt}\hrule height 0.3pt\vspace{3pt}
\textbf{Topic:} Computer Networks \\
\textbf{Subtopic:} \textit{Fragmentation and IP addressing, IPv4, CIDR notation}
\vspace{2pt}\hrule height 0.3pt\vspace{3pt}
\textbf{Guide / Hint:} \\
{\small At Router-1, max payload per fragment is $\lfloor (5000 - 20) / 8 \rfloor \times 8 = 4976$ bytes. 15000 bytes payload is fragmented into payloads of 4976, 4976, 4976, and 72 bytes (4 fragments). At Router-2 ($\text{MTU} = 3000$, max payload $\lfloor 2980/8 \rfloor \times 8 = 2976$ bytes), each of the three 4976-byte fragments splits into 2 fragments ($2976 + 2000$), yielding $3 \times 2 = 6$ fragments. The 72-byte fragment passes through intact. Total delivered fragments = $6 + 1 = 7$.}

\end{tcolorbox}
\vspace{6pt}\hrule height 0.35pt\vspace{6pt}
\noindent
\begin{tcolorbox}[
    blanker,
    sidebyside,
    sidebyside align=top,
    lefthand width=0.55\textwidth,
    top=4pt, bottom=4pt, left=0pt, right=0pt,
    middle=16pt,
    underlay={
        \draw[dashed, black!80, line width=0.8pt, dash pattern=on 4pt off 3pt]
            ([xshift=0.55\textwidth+8pt]frame.north west) --
            ([xshift=0.55\textwidth+8pt]frame.south west);
    }
]
\textbf{Q.58} \quad {\footnotesize\textsc{[NAT $\cdot$ 2 Mark(s) $\cdot$ No Neg]}} \hfill {\tiny\texttt{2025\_CS1\_CS\_Q58}}
\vspace{2pt}\hrule height 0.3pt\vspace{4pt}

Consider a probability distribution given by the density function $P(x)$:
\[ P(x) = \begin{cases} C x^2 & \text{for } 1 \le x \le 4 \\ 0 & \text{for } x < 1 \text{ or } x > 4 \end{cases} \]
The probability that $x$ lies between 2 and 3, i.e., $P(2 \le x \le 3)$ is \_\_\_\_\_\_\. (rounded off to three decimal places)

\par\vspace{6pt}\noindent\textbf{Your Answer:} \fbox{\rule{0pt}{1.3ex}\hspace{3cm}}

\tcblower

\textbf{Official Key:} \textbf{\large 0.300 to 0.302} \hfill \textit{\small [Medium]}
\vspace{2pt}\hrule height 0.3pt\vspace{3pt}
\textbf{Topic:} Probability and Statistics \\
\textbf{Subtopic:} \textit{Uniform, normal, exponential, Poisson and binomial distributions}
\vspace{2pt}\hrule height 0.3pt\vspace{3pt}
\textbf{Guide / Hint:} \\
{\small First find normalizing constant $C$: $\int_1^4 C x^2 \, dx = C \left[ \frac{x^3}{3} \right]_1^4 = C \frac{64 - 1}{3} = 21C = 1 \implies C = \frac{1}{21}$. Then $P(2 \le x \le 3) = \int_2^3 \frac{1}{21} x^2 \, dx = \frac{1}{21} \left[ \frac{x^3}{3} \right]_2^3 = \frac{27 - 8}{63} = \frac{19}{63} \approx 0.30158 \approx 0.302$.}

\end{tcolorbox}
\vspace{6pt}\hrule height 0.35pt\vspace{6pt}
\noindent
\begin{tcolorbox}[
    blanker,
    sidebyside,
    sidebyside align=top,
    lefthand width=0.55\textwidth,
    top=4pt, bottom=4pt, left=0pt, right=0pt,
    middle=16pt,
    underlay={
        \draw[dashed, black!80, line width=0.8pt, dash pattern=on 4pt off 3pt]
            ([xshift=0.55\textwidth+8pt]frame.north west) --
            ([xshift=0.55\textwidth+8pt]frame.south west);
    }
]
\textbf{Q.59} \quad {\footnotesize\textsc{[NAT $\cdot$ 2 Mark(s) $\cdot$ No Neg]}} \hfill {\tiny\texttt{2025\_CS1\_CS\_Q59}}
\vspace{2pt}\hrule height 0.3pt\vspace{4pt}

Consider a finite state machine (FSM) with one input $X$ and one output $f$, represented by the given state transition table:
\par\medskip
\begin{tabular}{|c|cc|cc|}
\hline
\textbf{Present state} & \multicolumn{2}{c|}{\textbf{Next state}} & \multicolumn{2}{c|}{\textbf{Output } $f$} \\
& $X=0$ & $X=1$ & $X=0$ & $X=1$ \\
\hline
A & F & B & 0 & 0 \\
B & D & C & 0 & 0 \\
C & F & E & 0 & 0 \\
D & G & A & 1 & 0 \\
E & D & C & 0 & 0 \\
F & F & B & 1 & 1 \\
G & G & H & 0 & 1 \\
H & G & A & 1 & 0 \\
\hline
\end{tabular}
\par\medskip
The minimum number of states required to realize this FSM is \_\_\_\_\_\_\. (Answer in integer)

\par\vspace{6pt}\noindent\textbf{Your Answer:} \fbox{\rule{0pt}{1.3ex}\hspace{3cm}}

\tcblower

\textbf{Official Key:} \textbf{\large 5 to 5} \hfill \textit{\small [Hard]}
\vspace{2pt}\hrule height 0.3pt\vspace{3pt}
\textbf{Topic:} Digital Logic \\
\textbf{Subtopic:} \textit{Combinational and sequential circuits}
\vspace{2pt}\hrule height 0.3pt\vspace{3pt}
\textbf{Guide / Hint:} \\
{\small Using state equivalence partitioning: initial partition by outputs groups states into $P_0 = \{(A, B, C, E), (D, H), (F), (G)\}$. Comparing transitions on $X=0, 1$: states $B$ and $E$ have identical transitions to $(D, C)$, while $A$ and $C$ transition to $(F, B)$ and $(F, E)$ (which are equivalent since $B \equiv E$). States $D$ and $H$ have identical next states $(G, A)$. Thus, equivalence classes are $\{A, C\}, \{B, E\}, \{D, H\}, \{F\}, \{G\}$, resulting in exactly 5 minimal states.}

\end{tcolorbox}
\vspace{6pt}\hrule height 0.35pt\vspace{6pt}
\noindent
\begin{tcolorbox}[
    blanker,
    sidebyside,
    sidebyside align=top,
    lefthand width=0.55\textwidth,
    top=4pt, bottom=4pt, left=0pt, right=0pt,
    middle=16pt,
    underlay={
        \draw[dashed, black!80, line width=0.8pt, dash pattern=on 4pt off 3pt]
            ([xshift=0.55\textwidth+8pt]frame.north west) --
            ([xshift=0.55\textwidth+8pt]frame.south west);
    }
]
\textbf{Q.60} \quad {\footnotesize\textsc{[NAT $\cdot$ 2 Mark(s) $\cdot$ No Neg]}} \hfill {\tiny\texttt{2025\_CS1\_CS\_Q60}}
\vspace{2pt}\hrule height 0.3pt\vspace{4pt}

Consider the given sequential circuit designed using D-Flip-flops. The circuit is initialized with some value (initial state). The number of distinct states the circuit will go through before returning back to the initial state is \_\_\_\_\_\_\. (Answer in integer)

\par\vspace{6pt}\noindent\textbf{Your Answer:} \fbox{\rule{0pt}{1.3ex}\hspace{3cm}}

\tcblower

\textbf{Official Key:} \textbf{\large 7 to 8} \hfill \textit{\small [Hard]}
\vspace{2pt}\hrule height 0.3pt\vspace{3pt}
\textbf{Topic:} Digital Logic \\
\textbf{Subtopic:} \textit{Combinational and sequential circuits}
\vspace{2pt}\hrule height 0.3pt\vspace{3pt}
\textbf{Guide / Hint:} \\
{\small Tracing the characteristic shift register feedback transitions where the inputs to the D flip-flops are determined by logic feedback of the current state $(Q_3, Q_2, Q_1, Q_0)$, the state cycle length forms a closed loop of period 7 (or 8 depending on the initial valid cycle).}

\end{tcolorbox}
\vspace{6pt}\hrule height 0.35pt\vspace{6pt}
\noindent
\begin{tcolorbox}[
    blanker,
    sidebyside,
    sidebyside align=top,
    lefthand width=0.55\textwidth,
    top=4pt, bottom=4pt, left=0pt, right=0pt,
    middle=16pt,
    underlay={
        \draw[dashed, black!80, line width=0.8pt, dash pattern=on 4pt off 3pt]
            ([xshift=0.55\textwidth+8pt]frame.north west) --
            ([xshift=0.55\textwidth+8pt]frame.south west);
    }
]
\textbf{Q.61} \quad {\footnotesize\textsc{[NAT $\cdot$ 2 Mark(s) $\cdot$ No Neg]}} \hfill {\tiny\texttt{2025\_CS1\_CS\_Q61}}
\vspace{2pt}\hrule height 0.3pt\vspace{4pt}

Consider the following C program:
\par\medskip
\begin{verbatim}
#include <stdio.h>
int foo(int S[], int size){
    if(size == 0) return 0;
    if(size == 1) return 1;
    if(S[0] != S[1]) return 1 + foo(S+1, size-1);
    return foo(S+1, size-1);
}
int main(){
    int A[] = {0, 1, 2, 2, 2, 0, 0, 1, 1};
    printf("%d", foo(A, 9));
    return 0;
}
\end{verbatim}
The value printed by the given C program is \_\_\_\_\_\_\. (Answer in integer)

\par\vspace{6pt}\noindent\textbf{Your Answer:} \fbox{\rule{0pt}{1.3ex}\hspace{3cm}}

\tcblower

\textbf{Official Key:} \textbf{\large 5 to 5} \hfill \textit{\small [Medium]}
\vspace{2pt}\hrule height 0.3pt\vspace{3pt}
\textbf{Topic:} Programming and Data Structures \\
\textbf{Subtopic:} \textit{Recursion}
\vspace{2pt}\hrule height 0.3pt\vspace{3pt}
\textbf{Guide / Hint:} \\
{\small The function `foo` counts 1 whenever two consecutive adjacent elements differ ($S[0] \neq S[1]$), plus 1 for the base case `size == 1`. In array $A = \{0, 1, 2, 2, 2, 0, 0, 1, 1\}$, adjacent transitions where $A[i] \neq A[i+1]$ occur at indices: $0 \to 1$ ($0 \neq 1$), $1 \to 2$ ($1 \neq 2$), $4 \to 5$ ($2 \neq 0$), and $6 \to 7$ ($0 \neq 1$). There are 4 such differences, plus 1 returned at `size == 1`, giving $4 + 1 = 5$.}

\end{tcolorbox}
\vspace{6pt}\hrule height 0.35pt\vspace{6pt}
\noindent
\begin{tcolorbox}[
    blanker,
    sidebyside,
    sidebyside align=top,
    lefthand width=0.55\textwidth,
    top=4pt, bottom=4pt, left=0pt, right=0pt,
    middle=16pt,
    underlay={
        \draw[dashed, black!80, line width=0.8pt, dash pattern=on 4pt off 3pt]
            ([xshift=0.55\textwidth+8pt]frame.north west) --
            ([xshift=0.55\textwidth+8pt]frame.south west);
    }
]
\textbf{Q.62} \quad {\footnotesize\textsc{[NAT $\cdot$ 2 Mark(s) $\cdot$ No Neg]}} \hfill {\tiny\texttt{2025\_CS1\_CS\_Q62}}
\vspace{2pt}\hrule height 0.3pt\vspace{4pt}

Let LIST be a datatype for an implementation of linked list defined as follows:
\begin{verbatim}
typedef struct list {
    int data;
    struct list *next;
} LIST;
\end{verbatim}
Suppose a program has created two linked lists, L1 and L2, whose contents are given (code for creating L1 and L2 is not provided here). L1 contains 9 nodes, and L2 contains 7 nodes.
Consider the following C program segment that modifies the list L1. The number of nodes that will be there in L1 after the execution of the code segment is \_\_\_\_\_\_\. (Answer in integer)
\par\medskip
\begin{verbatim}
int find (int query, LIST *list) {
    while (list != NULL){
        if(list->data == query) return 1;
        list = list->next;
    }
    return 0;
}
int main () {
    ...
    ptr1 = L1; ptr2 = L2;
    while (ptr1->next != NULL){
        query = ptr1->next->data;
        if (find (query, L2))
            ptr1->next = ptr1->next->next;
        else ptr1 = ptr1->next;
    }
    ...
    return 0;
}
\end{verbatim}

\par\vspace{6pt}\noindent\textbf{Your Answer:} \fbox{\rule{0pt}{1.3ex}\hspace{3cm}}

\tcblower

\textbf{Official Key:} \textbf{\large 5 to 5} \hfill \textit{\small [Medium]}
\vspace{2pt}\hrule height 0.3pt\vspace{3pt}
\textbf{Topic:} Programming and Data Structures \\
\textbf{Subtopic:} \textit{Arrays, stacks, queues, linked lists}
\vspace{2pt}\hrule height 0.3pt\vspace{3pt}
\textbf{Guide / Hint:} \\
{\small The loop inspects each node following the head of `L1` (`ptr1->next`). If its data value is found in `L2`, that node is deleted from `L1` by bypassing it (`ptr1->next = ptr1->next->next`). Among the 8 non-head nodes in `L1`, exactly 4 node values match elements in `L2` and are removed, leaving $9 - 4 = 5$ nodes remaining in `L1`.}

\end{tcolorbox}
\vspace{6pt}\hrule height 0.35pt\vspace{6pt}
\noindent
\begin{tcolorbox}[
    blanker,
    sidebyside,
    sidebyside align=top,
    lefthand width=0.55\textwidth,
    top=4pt, bottom=4pt, left=0pt, right=0pt,
    middle=16pt,
    underlay={
        \draw[dashed, black!80, line width=0.8pt, dash pattern=on 4pt off 3pt]
            ([xshift=0.55\textwidth+8pt]frame.north west) --
            ([xshift=0.55\textwidth+8pt]frame.south west);
    }
]
\textbf{Q.63} \quad {\footnotesize\textsc{[NAT $\cdot$ 2 Mark(s) $\cdot$ No Neg]}} \hfill {\tiny\texttt{2025\_CS1\_CS\_Q63}}
\vspace{2pt}\hrule height 0.3pt\vspace{4pt}

Consider the following C program:
\par\medskip
\begin{verbatim}
#include <stdio.h>
int gate (int n) {
    int d, t, newnum, turn;
    newnum = turn = 0; t = 1;
    while (n >= t) t *= 10;
    t /= 10;
    while (t > 0) {
        d = n / t;
        n = n % t;
        t /= 10;
        if (turn) newnum = 10 * newnum + d;
        turn = (turn + 1) % 2;
    }
    return newnum;
}
int main () {
    printf("%d", gate(14362));
    return 0;
}
\end{verbatim}
The value printed by the given C program is \_\_\_\_\_\_\. (Answer in integer)

\par\vspace{6pt}\noindent\textbf{Your Answer:} \fbox{\rule{0pt}{1.3ex}\hspace{3cm}}

\tcblower

\textbf{Official Key:} \textbf{\large 46 to 46} \hfill \textit{\small [Medium]}
\vspace{2pt}\hrule height 0.3pt\vspace{3pt}
\textbf{Topic:} Programming and Data Structures \\
\textbf{Subtopic:} \textit{Programming in C}
\vspace{2pt}\hrule height 0.3pt\vspace{3pt}
\textbf{Guide / Hint:} \\
{\small For $n = 14362$, the first loop computes $t = 10000$. The second loop extracts digits from left to right: $d=1$ (`turn=0`, toggles to 1), $d=4$ (`turn=1`, `newnum = 4`, toggles to 0), $d=3$ (`turn=0`, toggles to 1), $d=6$ (`turn=1`, `newnum = 40 + 6 = 46`, toggles to 0), $d=2$ (`turn=0`, toggles to 1). The function extracts even-indexed digits (1-based: 2nd and 4th digits), producing 46.}

\end{tcolorbox}
\vspace{6pt}\hrule height 0.35pt\vspace{6pt}
\noindent
\begin{tcolorbox}[
    blanker,
    sidebyside,
    sidebyside align=top,
    lefthand width=0.55\textwidth,
    top=4pt, bottom=4pt, left=0pt, right=0pt,
    middle=16pt,
    underlay={
        \draw[dashed, black!80, line width=0.8pt, dash pattern=on 4pt off 3pt]
            ([xshift=0.55\textwidth+8pt]frame.north west) --
            ([xshift=0.55\textwidth+8pt]frame.south west);
    }
]
\textbf{Q.64} \quad {\footnotesize\textsc{[NAT $\cdot$ 2 Mark(s) $\cdot$ No Neg]}} \hfill {\tiny\texttt{2025\_CS1\_CS\_Q64}}
\vspace{2pt}\hrule height 0.3pt\vspace{4pt}

The maximum value of $x$ such that the edge between the nodes B and C is included in every minimum spanning tree of the given graph is \_\_\_\_\_\_\. (answer in integer)

\par\vspace{6pt}\noindent\textbf{Your Answer:} \fbox{\rule{0pt}{1.3ex}\hspace{3cm}}

\tcblower

\textbf{Official Key:} \textbf{\large 5 to 5} \hfill \textit{\small [Medium]}
\vspace{2pt}\hrule height 0.3pt\vspace{3pt}
\textbf{Topic:} Algorithms \\
\textbf{Subtopic:} \textit{Graph traversals, minimum spanning trees, shortest paths}
\vspace{2pt}\hrule height 0.3pt\vspace{3pt}
\textbf{Guide / Hint:} \\
{\small By the cycle property of minimum spanning trees, an edge $e = (B, C)$ is included in every MST if and only if it is strictly lighter than the heaviest edge on every cycle containing $(B, C)$. Evaluating the fundamental cycle formed by adding $(B, C)$ where the alternative path has maximum edge weight 6, we require $x < 6$, so the maximum integer value for $x$ is 5.}

\end{tcolorbox}
\vspace{6pt}\hrule height 0.35pt\vspace{6pt}
\noindent
\begin{tcolorbox}[
    blanker,
    sidebyside,
    sidebyside align=top,
    lefthand width=0.55\textwidth,
    top=4pt, bottom=4pt, left=0pt, right=0pt,
    middle=16pt,
    underlay={
        \draw[dashed, black!80, line width=0.8pt, dash pattern=on 4pt off 3pt]
            ([xshift=0.55\textwidth+8pt]frame.north west) --
            ([xshift=0.55\textwidth+8pt]frame.south west);
    }
]
\textbf{Q.65} \quad {\footnotesize\textsc{[NAT $\cdot$ 2 Mark(s) $\cdot$ No Neg]}} \hfill {\tiny\texttt{2025\_CS1\_CS\_Q65}}
\vspace{2pt}\hrule height 0.3pt\vspace{4pt}

In a double hashing scheme, $h_1(k) = k \bmod 11$ and $h_2(k) = 1 + (k \bmod 7)$ are the auxiliary hash functions. The size $m$ of the hash table is 11. The hash function for the $i$-th probe in the open address table is $[h_1(k) + i \cdot h_2(k)] \bmod m$. The following keys are inserted in the given order: 63, 50, 25, 79, 67, 24.
The slot at which key 24 gets stored is \_\_\_\_\_\_\. (Answer in integer)

\par\vspace{6pt}\noindent\textbf{Your Answer:} \fbox{\rule{0pt}{1.3ex}\hspace{3cm}}

\tcblower

\textbf{Official Key:} \textbf{\large 10 to 10} \hfill \textit{\small [Medium]}
\vspace{2pt}\hrule height 0.3pt\vspace{3pt}
\textbf{Topic:} Algorithms \\
\textbf{Subtopic:} \textit{Searching, sorting, hashing}
\vspace{2pt}\hrule height 0.3pt\vspace{3pt}
\textbf{Guide / Hint:} \\
{\small Hash values: key 63 at $63 \bmod 11 = 8$; key 50 at $50 \bmod 11 = 6$; key 25 at $25 \bmod 11 = 3$; key 79 at $79 \bmod 11 = 2$; key 67 has $h_1(67) = 1$, stores at 1. For key 24: $h_1(24) = 24 \bmod 11 = 2$ (collision with 79), $h_2(24) = 1 + (24 \bmod 7) = 1 + 3 = 4$. Probe $i=1$: $(2 + 4) \bmod 11 = 6$ (collision with 50). Probe $i=2$: $(2 + 8) \bmod 11 = 10$ (empty). Hence, key 24 is stored in slot 10.}

\end{tcolorbox}
\vspace{6pt}\hrule height 0.35pt\vspace{6pt}
\end{document}